\section{Critères différentiels de lissité formelle et de régularité}
\label{section:0.22-fr}

Les principaux résultats de ce paragraphe sont~:
\begin{enumerate}
  \item[a)] Les critères (22.2.2) et (22.5.8), qui donnent des conditions
  nécessaires et suffisantes pour qu'un anneau local noethérien complet $A$
  contenant un corps $k$ soit formellement lisse sur $k$~; le critère
  (22.5.8) s'énonce sans faire intervenir de notions différentielles, et
  complète (19.6.6)~; l'énoncé (22.2.2), de nature un peu plus technique,
  permet de donner une \emph{classification}, pour $k$ fixé, des anneaux
  locaux $A$ considérés (22.2.5)~: ils sont donnés par leur dimension $n$ et
  leur corps résiduel $K$, sous la seule condition
  $n\geqslant\operatorname{rg}_K\Upsilon_{K/k}$.

  \item[b)] Le théorème (22.3.3), et le corollaire (22.7.6) du critère
  jacobien de Nagata (22.7.3), qui donnent des conditions permettant
  d'affirmer qu'un anneau \emph{localisé} d'un anneau local noethérien complet
  $A$ contenant un corps $k$ est géométriquement régulier sur $k$~; ils
  joueront un rôle important dans l'étude des propriétés «~fines~» des
  anneaux locaux faite au chap.~IV, §~7. On ne possède pour le moment aucune
  démonstration de (22.7.6) indépendante du critère jacobien de Nagata.

  \item[c)] Le critère jacobien de Zariski (22.6.7) et ses variantes, qui
  sont des conséquences faciles de (20.7.8).
\end{enumerate}

\oldpage[0\textsubscript{IV-1}]{183}

\subsection{Relèvement de la lissité formelle}
\label{subsection:0.22.1-fr}

\begin{theorem}[22.1.1]
\label{0.22.1.1-fr}
Soient $s:\Lambda\to A$, $u:A\to B$ deux homomorphismes d'anneaux,
$\mathfrak J$ un idéal de $B$, $C$ l'anneau quotient $B/\mathfrak J$. On
suppose que $C$ est une $\Lambda$-algèbre formellement lisse pour la topologie
discrète.
\begin{enumerate}
  \item[(i)] Supposons vérifiées les conditions suivantes~:
  \begin{enumerate}
    \item[$\alpha$)] $\mathfrak J/\mathfrak J^2$ est un $C$-module projectif.
    \item[$\beta$)] L'homomorphisme canonique
    $\mathbf S_C^*(\mathfrak J/\mathfrak J^2)\to
    \operatorname{gr}_{\mathfrak J}^*(B)$ (19.5.2) est bijectif.
    \item[$\gamma$)] L'homomorphisme caractéristique
    $\chi:\Upsilon_{C/B/\Lambda}\to\mathfrak J/\mathfrak J^2$ (20.6.20)
    est injectif.
    \item[$\delta$)] Le $C$-module
    $\Omega^1_{B/A}\otimes_BC$ est projectif.
  \end{enumerate}
  Alors $B$, muni de la topologie $\mathfrak J$-préadique, est une
  $\Lambda$-algèbre formellement lisse.

  \item[(ii)] Inversement, supposons que $B$ soit une $\Lambda$-algèbre
  formellement lisse pour la topologie $\mathfrak J$-préadique, et que $A$
  soit une $\Lambda$-algèbre formellement lisse pour la topologie discrète.
  Alors les conditions $\alpha$), $\beta$), $\gamma$), $\delta$) de (i) sont
  vérifiées.
\end{enumerate}
\end{theorem}

\begin{proof}
En vertu de la suite exacte (20.6.22.1) (qui est applicable, puisque
$B/\mathfrak J^2$ est une extension $\Lambda$-triviale de $C$ par
$\mathfrak J/\mathfrak J^2$, $C$ étant supposée être une $\Lambda$-algèbre
formellement lisse (19.4.4)), la condition $\gamma$) équivaut à
$\Upsilon^C_{B/A/\Lambda}=0$, ou encore à la suivante~:
\begin{enumerate}
  \item[$\gamma'$)] L'homomorphisme
  \[
    u_{B/A/\Lambda}\otimes 1_C:
    \Omega^1_{A/\Lambda}\otimes_A C
    \longrightarrow \Omega^1_{B/\Lambda}\otimes_B C
  \]
  est injectif.
\end{enumerate}

(i) Les conditions $\alpha$) et $\beta$) entraînent que $B$ est une
$\Lambda$-algèbre formellement lisse pour la topologie
$\mathfrak J$-préadique (19.5.4). Il revient donc au même de dire que $B$ est
(pour la topologie $\mathfrak J$-préadique) une $A$-algèbre formellement
lisse, ou une $A$-algèbre formellement lisse \emph{relativement à $\Lambda$}.
Pour voir que $B$ est une $A$-algèbre formellement lisse, il suffit donc, en
vertu de (20.7.2), de prouver que l'homomorphisme continu de $B$-modules
topologiques
\[
  u_{B/A/\Lambda}:\Omega^1_{A/\Lambda}\otimes_A B
  \longrightarrow\Omega^1_{B/\Lambda}
\]
est \emph{formellement inversible à gauche}. Or, puisque $B$ est une
$\Lambda$-algèbre formellement lisse, le $B$-module topologique
$\Omega^1_{B/\Lambda}$ est formellement projectif (20.4.9) et sa topologie
est déduite de celle de $B$ (20.4.5). En vertu de (19.1.9), il suffit donc,
pour que $u_{B/A/\Lambda}$ soit formellement inversible à gauche, que
$u_{B/A/\Lambda}\otimes 1_C$ soit inversible à gauche. Mais en vertu de
$\gamma'$), cette dernière application est injective, donc on a la suite
exacte (20.6.14.7)
\[
  0\longrightarrow\Omega^1_{A/\Lambda}\otimes_A C
  \longrightarrow\Omega^1_{B/\Lambda}\otimes_B C
  \longrightarrow\Omega^1_{B/A}\otimes_B C\longrightarrow0.
\]
Enfin, en vertu de $\delta$), le $C$-module
$\Omega^1_{B/A}\otimes_BC$ est projectif, donc la suite exacte précédente
est scindée, ce qui achève de prouver (i).

\oldpage[0\textsubscript{IV-1}]{184}

(ii) Si $B$ est une $\Lambda$-algèbre formellement lisse pour la topologie
$\mathfrak J$-préadique, et $A$ une $\Lambda$-algèbre formellement lisse pour
la topologie discrète, $B$ est une $A$-algèbre formellement lisse pour la
topologie $\mathfrak J$-préadique (19.3.5, (ii))~; les conditions $\alpha$) et
$\beta$) résultent donc de (19.5.6). En outre, (20.7.2) montre que
$u_{B/A/\Lambda}$ est formellement inversible à gauche, donc
$u_{B/A/\Lambda}\otimes 1_C$ est inversible à gauche (19.1.7), et \emph{a
fortiori} injectif. Enfin, $\Omega^1_{B/A}$ est un $B$-module formellement
projectif (20.4.9), et par suite $\Omega^1_{B/A}\otimes_BC$ est un $C$-module
projectif (19.2.4).
\end{proof}

\begin{corollary}[22.1.2]
\label{0.22.1.2-fr}
Soient $s:\Lambda\to A$, $u:A\to B$ deux homomorphismes d'anneaux,
$\mathfrak m$ un idéal maximal de $B$, $K=B/\mathfrak m$ le corps quotient.
On suppose que les $\Lambda$-algèbres $A$ et $K$ soient formellement lisses
pour la topologie discrète. Alors les conditions suivantes sont
équivalentes~:
\begin{enumerate}
  \item[a)] $B$ est une $A$-algèbre formellement lisse pour la topologie
  $\mathfrak m$-préadique.
  \item[b)] L'homomorphisme canonique
  \begin{equation}
  \tag{22.1.2.1}
  \label{0.22.1.2.1-fr}
    \mathbf S_K^*(\mathfrak m/\mathfrak m^2)
    \longrightarrow\operatorname{gr}_{\mathfrak m}^*(B)
  \end{equation}
  est bijectif, et l'homomorphisme caractéristique
  \begin{equation}
  \tag{22.1.2.2}
  \label{0.22.1.2.2-fr}
    \chi:\Upsilon_{K/B/\Lambda}\longrightarrow\mathfrak m/\mathfrak m^2
  \end{equation}
  est injectif.
\end{enumerate}
\end{corollary}

\begin{proof}
Cela résulte aussitôt de (22.1.1), les conditions $\alpha$) et $\delta$)
étant ici trivialement vérifiées, puisque $K$ est un corps.
\end{proof}

\begin{remark}[22.1.3]
\label{0.22.1.3-fr}
Supposons vérifiées les hypothèses de (22.1.2) et en outre supposons que
$\mathfrak m/\mathfrak m^2$ soit un $K$-espace vectoriel de rang \emph{fini}.
Alors, pour que $B$ soit une $A$-algèbre formellement lisse pour la topologie
$\mathfrak m$-préadique, il suffit que les conditions suivantes soient
satisfaites~:

\begin{env}[22.1.3.1]
\label{0.22.1.3.1-fr}
Étant donnés une algèbre de polynômes $E=K[T_1,\ldots,T_n]$, l'idéal maximal
$\mathfrak n$ de $E$ engendré par les $T_i$ $(1\leqslant i\leqslant n)$, et
un idéal $\mathfrak b\supset\mathfrak n^k$ dans $E$, tout
$\Lambda$-homomorphisme $v:B\to E/\mathfrak b$ qui, par passage aux
quotients, donne l'identité $K\to K$, se factorise en
\[
  B\xrightarrow{w}E/\mathfrak n^k\longrightarrow E/\mathfrak b,
\]
où $w$ est un $\Lambda$-homomorphisme.
\end{env}

\begin{env}[22.1.3.2]
\label{0.22.1.3.2-fr}
Si $F$ est l'extension triviale $D_K(K)$ («~algèbre des nombres duaux sur
$K$~»), $\rho:A\to F$ un homomorphisme d'anneaux tel que le composé
$A\xrightarrow{\rho}F\to K$ soit l'homomorphisme canonique, alors tout
$\Lambda$-homomorphisme $v:B\to K$ qui, par passage aux quotients, donne
l'identité $K\to K$, se factorise en
\[
  B\xrightarrow{w}F\longrightarrow K,
\]
où $w$ est un $\Lambda$-homomorphisme (pour la structure de $A$-algèbre sur
$F$ définie par $\rho$).
\end{env}

En effet, on a vu (19.5.5) que la condition (22.1.3.1) entraîne que
l'homomorphisme canonique (22.1.2.1) est bijectif. En vertu de (22.1.2), il
suffit ensuite de voir que l'homomorphisme canonique
\[
  u_{B/A/\Lambda}\otimes1_K:
  \Omega^1_{A/\Lambda}\otimes_AK
  \longrightarrow\Omega^1_{B/\Lambda}\otimes_BK
\]
est injectif, ou encore (puisqu'il s'agit d'espaces vectoriels sur $K$) que
pour tout $K$-espace vectoriel $L$, l'homomorphisme canonique
$\operatorname{D\acute{e}r}_\Lambda(B,L)\to\operatorname{D\acute{e}r}_\Lambda(A,L)$ est
\emph{surjectif} (en vertu de (20.4.8))~; il est clair que pour cela, il faut
et il suffit que l'homomorphisme canonique
$\operatorname{D\acute{e}r}_\Lambda(B,K)\to\operatorname{D\acute{e}r}_\Lambda(A,K)$ soit
surjectif, d'où notre assertion, en vertu de (20.1.5).

On notera que les deux conditions (22.1.3.1) et (22.1.3.2) sont entraînées
par la suivante~:

\begin{env}[22.1.3.3]
\label{0.22.1.3.3-fr}
Pour toute $\Lambda$-algèbre locale artinienne $E$ de corps résiduel égal à
$K$, et tout idéal nilpotent $\mathfrak r$ de $E$, tout
$\Lambda$-homomorphisme $B\to E/\mathfrak r$ qui, par passage aux quotients,
donne l'identité $K\to K$, se factorise en
$B\xrightarrow{u}E\to E/\mathfrak r$, où $u$ est un
$\Lambda$-homomorphisme.
\end{env}
\end{remark}

\oldpage[0\textsubscript{IV-1}]{185}

Ce résultat se généralise comme suit~:

\begin{proposition}[22.1.4]
\label{0.22.1.4-fr}
Soit $\varphi:A\to B$ un homomorphisme local d'anneaux locaux noethériens,
$k$ (resp. $K$) le corps résiduel de $A$ (resp. $B$). Pour que $B$ soit une
$A$-algèbre formellement lisse (pour les topologies préadiques), il faut et
il suffit que pour tout anneau local artinien $C$, de corps résiduel égal à
$K$, tout homomorphisme local $A\to C$ (faisant de $C$ une $A$-algèbre
topologique), et tout idéal nilpotent $\mathfrak J$ de $C$, tout
$A$-homomorphisme local $B\to C/\mathfrak J$ tel que l'homomorphisme
$K\to K$ obtenu par passage aux quotients soit l'identité, se factorise en
\[
  B\xrightarrow{f}C\longrightarrow C/\mathfrak J,
\]
où $f$ est un $A$-homomorphisme local.
\end{proposition}

\begin{proof}
La condition étant nécessaire par définition (19.3.1), montrons qu'elle est
suffisante. Comme les $\widehat A$-homomorphismes locaux de $\widehat B$ dans
une $\widehat A$-algèbre locale discrète proviennent par extension des
$A$-homomorphismes locaux de $B$ dans cette algèbre, on peut se borner au cas
où $A$ et $B$ sont \emph{complets}, en vertu de (19.3.6). Lorsque $A$ est un
\emph{corps} $k$, la proposition résulte de (22.1.3.3), où l'on prend pour
$\Lambda$ le sous-corps premier de $k$, et de (19.6.1).

Dans le cas général, posons $B_0=B\otimes_Ak=B/\mathfrak mB$
($\mathfrak m$ idéal maximal de $A$), qui est complet. Si $C_0$ est une
$k$-algèbre locale artinienne de corps résiduel égal à $K$, $\mathfrak J_0$
un idéal nilpotent de $C_0$, tout $k$-homomorphisme local
$g_0:B_0\to C_0/\mathfrak J_0$ donnant l'identité sur $K$ par passage aux
quotients, donne par composition un $A$-homomorphisme local
\[
  g:B\longrightarrow B_0\xrightarrow{g_0}C_0/\mathfrak J_0
\]
qui par hypothèse se factorise donc en
$B\xrightarrow{f}C_0\to C_0/\mathfrak J_0$, où $f$ est un
$A$-homomorphisme local~; mais comme $C_0$ est une $k$-algèbre, $f$ se
factorise en $B\to B_0\xrightarrow{f_0}C_0$, où $f_0$ est un
$k$-homomorphisme local. On voit donc que les hypothèses de l'énoncé sont
encore remplies lorsqu'on y substitue $B_0$ et $k$ à $B$ et $A$
respectivement~; donc $B_0$ est une $k$-algèbre \emph{formellement lisse}
d'après ce que l'on vient de voir.

Appliquant (19.7.2) et (19.7.1), on voit qu'il existe un anneau local
noethérien complet $B'$, un homomorphisme local $A\to B'$ faisant de $B'$ un
$A$-module \emph{plat} et une $A$-algèbre \emph{formellement lisse}, et un
$k$-isomorphisme
\[
  u'_0:B'_0=B'\otimes_Ak\xrightarrow{\sim}B_0.
\]
Soit $v'_0$ l'isomorphisme réciproque de $u'_0$~; nous nous proposons de
montrer qu'il existe un $A$-homomorphisme local $v:B\to B'$ tel que $v'_0$
provienne de $v$ par passage aux quotients. Pour cela, désignons par
$\mathfrak n$ et $\mathfrak n'$ les idéaux maximaux respectifs de $B$ et
$B'$, et, pour tout $j$, soit
\[
  w_j:B/(\mathfrak n^{j+1}+\mathfrak mB)
  \longrightarrow B'/(\mathfrak n'^{j+1}+\mathfrak mB')
\]
l'homomorphisme déduit de $v'_0$ par passage aux quotients~; il suffira de
former pour tout $j$ un $A$-homomorphisme local
$v_j:B/\mathfrak n^{j+1}\to B'/\mathfrak n'^{j+1}$ tel que les diagrammes
\[
\xymatrix@C=4em@R=3em{
  B/\mathfrak n^{j+1} \ar[r]^{v_j} \ar[d] &
  B'/\mathfrak n'^{j+1} \ar[d] \\
  B/\mathfrak n^j \ar[r]_{v_{j-1}} & B'/\mathfrak n'^j
}
\qquad
\xymatrix@C=4em@R=3em{
  B/\mathfrak n^{j+1} \ar[r]^{v_j} \ar[d] &
  B'/\mathfrak n'^{j+1} \ar[d] \\
  B/(\mathfrak n^{j+1}+\mathfrak mB) \ar[r]_{w_j} &
  B'/(\mathfrak n'^{j+1}+\mathfrak mB')
}
\]
soient commutatifs~; $B$ et $B'$ étant complets, l'homomorphisme
$v=\varprojlim v_j$ répondra à la question.

\oldpage[0\textsubscript{IV-1}]{186}

Or, la formation par récurrence de $v_j$ résulte de l'hypothèse sur $B$ et
du même raisonnement que dans (19.3.10.1), en utilisant le lemme
(19.3.10.2) et le fait que
\[
  \mathfrak n'^j+(\mathfrak n'^{j+1}+\mathfrak mB')
  =\mathfrak n'^j+\mathfrak mB'.
\]
Il résulte alors de (19.7.1.4) que $v$ est un $A$-\emph{isomorphisme}, car
$B'$ est un $A$-module plat, et $B$ est complet pour la topologie
$\mathfrak m$-préadique, les idéaux $\mathfrak m^iB$ étant fermés dans $B$
pour la topologie $\mathfrak n$-adique ($\mathbf 0_{\mathrm I}$, 7.3.5).
C.Q.F.D.
\end{proof}

\subsection{Caractérisation différentielle des algèbres locales formellement
lisses sur un corps}
\label{subsection:0.22.2-fr}

\begin{env}[22.2.1]
\label{0.22.2.1-fr}
Soient $k$ un corps, $P$ son sous-corps premier, $A$ une $k$-algèbre qui est
un anneau \emph{local}, $\mathfrak m$ son idéal maximal, $K=A/\mathfrak m$
son corps résiduel. Comme $K$ est séparable sur $P$, $K$ est une
$P$-algèbre formellement lisse pour la topologie discrète (19.6.1), et par
suite la $P$-extension $A/\mathfrak m^2$ de $K$ par
$\mathfrak m/\mathfrak m^2$ est $P$-\emph{triviale}, et l'homomorphisme
caractéristique
\begin{equation}
\tag{22.2.1.1}
\label{0.22.2.1.1-fr}
  \chi_{A/k}:\Upsilon_{K/k}\longrightarrow\mathfrak m/\mathfrak m^2
\end{equation}
est donc défini (20.6.23).
\end{env}

\begin{theorem}[22.2.2]
\label{0.22.2.2-fr}
Sous les conditions de (22.2.1), pour que $A$ soit une $k$-algèbre
formellement lisse pour la topologie $\mathfrak m$-préadique, il faut et il
suffit que les conditions suivantes soient satisfaites~:
\begin{enumerate}
  \item[(i)] L'homomorphisme canonique
  $\mathbf S_K^*(\mathfrak m/\mathfrak m^2)\to
  \operatorname{gr}_{\mathfrak m}^*(A)$ est bijectif.
  \item[(ii)] L'homomorphisme caractéristique
  $\chi_{A/k}:\Upsilon_{K/k}\to\mathfrak m/\mathfrak m^2$ est injectif.
\end{enumerate}
\end{theorem}

\begin{proof}
C'est un cas particulier de (22.1.2), où l'on remplace $\Lambda$ par $P$,
$A$ par $k$ et $B$ par $A$~; $K$ et $k$, étant séparables sur $P$, sont en
effet des $P$-algèbres formellement lisses (19.6.1).
\end{proof}

\begin{remark}[22.2.3]
\label{0.22.2.3-fr}
La condition (ii) de (22.2.2) peut être remplacée par une quelconque des
suivantes~:
\begin{enumerate}
  \item[a)] L'homomorphisme canonique
  $\Omega^1_k\otimes_kK\to\Omega^1_A\otimes_AK$ est injectif.
  \item[b)] Pour tout $n$, l'homomorphisme canonique
  \[
    \Omega^1_k\otimes_k(A/\mathfrak m^{n+1})
    \longrightarrow
    \Omega^1_A\otimes_A(A/\mathfrak m^{n+1})
  \]
  est inversible à gauche.
  \item[c)] L'homomorphisme canonique
  $\Omega^1_k\widehat\otimes_kA\to\widehat\Omega^1_A$ est inversible à
  gauche.
  \item[d)] $A$ est une $k$-algèbre formellement lisse (pour la topologie
  $\mathfrak m$-préadique) \emph{relativement au corps premier}.
\end{enumerate}
En effet, on a déjà remarqué dans (22.1.1) que l'injectivité de $\chi_{A/k}$
est équivalente à $a)$, et que la conjonction de $a)$ et de la condition (i)
de (22.2.2) équivaut à celle de $d)$ et de (i). Enfin, on a aussi vu dans
(22.1.1) que $b)$ et $d)$ sont équivalentes, et l'équivalence de $b)$ et
$c)$ résulte de (19.1.8).
\end{remark}
\begin{env}[22.2.4]
\label{0.22.2.4-fr}
La principale application de (22.2.2) a trait aux $k$-algèbres locales
noethériennes et \emph{complètes} (cas où $\mathfrak m/\mathfrak m^2$ est un
$K$-espace vectoriel de rang \emph{fini}). De façon précise, soient $k$ un
corps, $K$ une extension de $k$, $V$ un $K$-espace vectoriel de

\oldpage[0\textsubscript{IV-1}]{187}

rang fini~; nous considérons les triplets $(A,\alpha,\beta)$, où $A$ est une
$k$-algèbre locale noethérienne complète \emph{formellement lisse} d'idéal
maximal $\mathfrak m$, $\alpha:A/\mathfrak m\xrightarrow{\sim}K$ un
isomorphisme de $k$-algèbres,
$\beta:\mathfrak m/\mathfrak m^2\xrightarrow{\sim}V$ un di-isomorphisme
d'espaces vectoriels (pour l'isomorphisme $\alpha$). Étant donnés deux tels
triplets $(A,\alpha,\beta)$, $(A',\alpha',\beta')$, une \emph{équivalence} du
premier sur le second est un $k$-isomorphisme $u:A\xrightarrow{\sim}A'$ tel
que (si $\mathfrak m'$ est l'idéal maximal de $A'$), les diagrammes
\[
\xymatrix@C=4em@R=3em{
  A/\mathfrak m \ar[r]^{\operatorname{gr}^0(u)} \ar[d]_{\alpha} &
  A'/\mathfrak m' \ar[d]^{\alpha'} \\
  K \ar[r]_{1_K} & K
}
\qquad
\xymatrix@C=4em@R=3em{
  \mathfrak m/\mathfrak m^2 \ar[r]^{\operatorname{gr}^1(u)}
    \ar[d]_{\beta} &
  \mathfrak m'/\mathfrak m'^2 \ar[d]^{\beta'} \\
  V \ar[r]_{1_V} & V
}
\]
sont commutatifs. Il est immédiat que les \emph{classes d'équivalence} des
triplets $(A,\alpha,\beta)$ forment un \emph{ensemble} que nous noterons
$\operatorname{FL}(K/k,V)$. Remarquons maintenant qu'à tout triplet
$(A,\alpha,\beta)$ est associé le $K$-homomorphisme composé d'espaces
vectoriels
\[
  \Upsilon_{K/k}\xrightarrow{\chi_{A/k}}
  \mathfrak m/\mathfrak m^2\xrightarrow{\beta}V
\]
(où $\chi_{A/k}$ est considéré comme un di-homomorphisme relatif à
$\alpha^{-1}$)~; la définition précédente prouve (par (20.6.22.2)) que ce
$K$-homomorphisme \emph{ne dépend que de la classe d'équivalence} de
$(A,\alpha,\beta)$ dans $\operatorname{FL}(K/k,V)$~; autrement dit on a
défini une application canonique
\begin{equation}
\tag{22.2.4.1}
\label{0.22.2.4.1-fr}
  \operatorname{FL}(K/k,V)\longrightarrow
  \operatorname{Hom}_K(\Upsilon_{K/k},V).
\end{equation}
\end{env}

\begin{proposition}[22.2.5]
\label{0.22.2.5-fr}
Soient $k$ un corps, $K$ une extension de $k$, $V$ un $K$-espace vectoriel
de rang fini. L'application canonique (22.2.4.1) est une bijection de
$\operatorname{FL}(K/k,V)$ sur l'ensemble des $K$-homomorphismes injectifs
de $\Upsilon_{K/k}$ dans $V$.
\end{proposition}

\begin{proof}
(i) Pour montrer que (22.2.4.1) est injective, il faut prouver ce qui suit~:
soient $A$, $A'$ deux $k$-algèbres locales noethériennes complètes
formellement lisses, $\mathfrak m$, $\mathfrak m'$ leurs idéaux maximaux
respectifs, $\alpha:A/\mathfrak m\xrightarrow{\sim}K$,
$\alpha':A'/\mathfrak m'\xrightarrow{\sim}K$ deux $k$-isomorphismes~;
supposons donnés deux $k$-isomorphismes
\[
  v^0:A/\mathfrak m\xrightarrow{\sim}A'/\mathfrak m',\qquad
  v^1:\mathfrak m/\mathfrak m^2\xrightarrow{\sim}
      \mathfrak m'/\mathfrak m'^2
\]
tels que les diagrammes
\begin{equation}
\tag{22.2.5.1}
\label{0.22.2.5.1-fr}
\xymatrix@C=4em@R=3em{
  A/\mathfrak m \ar[r]^{v^0} \ar[dr]_{\alpha} &
  A'/\mathfrak m' \ar[dl]^{\alpha'} \\
  & K
}
\qquad
\xymatrix@C=4em@R=3em{
  \mathfrak m/\mathfrak m^2 \ar[r]^{v^1} &
  \mathfrak m'/\mathfrak m'^2 \\
  \Upsilon_{K/k} \ar[u]^{\chi_A} \ar[ur]_{\chi_{A'}} &
}
\end{equation}
soient commutatifs. Il s'agit de prouver qu'il existe un $k$-isomorphisme
$u:A\xrightarrow{\sim}A'$ tel que $\operatorname{gr}^0(u)=v^0$ et
$\operatorname{gr}^1(u)=v^1$. Remarquons en premier lieu que puisque $K$ est
un corps, l'homomorphisme canonique (20.6.7.1) est \emph{bijectif} (20.6.11),
donc il existe déjà un $k$-isomorphisme
$v:A/\mathfrak m^2\xrightarrow{\sim}A'/\mathfrak m'^2$ tel que
$\operatorname{gr}^0(v)=v^0$ et $\operatorname{gr}^1(v)=v^1$. Notons
maintenant que puisque $A'$ est métrisable et complète, et $A$ une
$k$-algèbre formellement lisse, l'homomorphisme composé
$A\to A/\mathfrak m^2\xrightarrow{v}A'/\mathfrak m'^2$ se factorise en
$A\xrightarrow{u}A'\to A'/\mathfrak m'^2$, où $u$ est

\oldpage[0\textsubscript{IV-1}]{188}

un $k$-homomorphisme (19.3.11), et l'on a donc déjà
$\operatorname{gr}^0(u)=v^0$ et $\operatorname{gr}^1(u)=v^1$. Or, on a le
diagramme commutatif
\[
\xymatrix@C=5em@R=3em{
  \mathbf S_K^*(\mathfrak m/\mathfrak m^2) \ar[r] \ar[d]_{w} &
  \operatorname{gr}_{\mathfrak m}^*(A) \ar[d]^{\operatorname{gr}^*(u)} \\
  \mathbf S_K^*(\mathfrak m'/\mathfrak m'^2) \ar[r] &
  \operatorname{gr}_{\mathfrak m'}^*(A')
}
\]
où les flèches horizontales sont les homomorphismes canoniques, et $w$
provient de $v^0$ et $v^1$~; comme $v^0$ et $v^1$ sont bijectifs, il en est
de même de $w$, et puisque $A$ et $A'$ sont des $k$-algèbres formellement
lisses, on déduit de (22.2.2, (i)) que $\operatorname{gr}(u)$ est bijectif~;
mais comme $A$ et $A'$ sont séparés et complets, cela entraîne que $u$
lui-même est \emph{bijectif} (Bourbaki, \emph{Alg. comm.}, chap.~III, §~2,
n\textsuperscript{o}~8, cor.~3 du th.~1).

(ii) Reste à montrer que l'image de (22.2.4.1) est l'ensemble des
homomorphismes \emph{injectifs} de $\Upsilon_{K/k}$ dans $V$~; on sait déjà
par (22.2.2, (ii)) qu'elle est contenue dans cet ensemble~; il suffira alors
de prouver le lemme suivant.
\end{proof}

\begin{lemma}[22.2.5.2]
\label{0.22.2.5.2-fr}
Pour tout $K$-homomorphisme $h:\Upsilon_{K/k}\to V$, il existe une
$k$-algèbre $A$ qui est un anneau local noethérien régulier et complet,
d'idéal maximal $\mathfrak m$, de corps résiduel $K$, et un $K$-isomorphisme
$\beta:\mathfrak m/\mathfrak m^2\xrightarrow{\sim}V$ tels que
$h=\beta\circ\chi_{A/k}$.
\end{lemma}

\begin{proof}
En effet, si ce lemme est prouvé, le fait que $A$ est régulier entraîne que
la condition (i) de (22.2.2) est satisfaite (17.1.1)~; si en outre $h$ est
injectif, (22.2.2) montrera que $A$ est une $k$-algèbre formellement lisse
dont la classe aura $h$ pour image (par contre si $h$ n'est pas injectif, la
$k$-algèbre $A$ dont (22.2.5.2) prouve l'existence n'est pas formellement
lisse).

Pour prouver (22.2.5.2), soient $B$ la $K$-algèbre $\mathbf S_K^*(V)$,
$\mathfrak n$ l'idéal d'augmentation de $B$ et $A=\widehat B$ la
$K$-algèbre complétée de $B$ pour la topologie $\mathfrak n$-préadique (de
sorte que $A$ est isomorphe à l'algèbre de séries formelles
$K[[T_1,\ldots,T_n]]$ si $n=\operatorname{rg}_KV$). En vertu de (20.6.11)
il existe sur $A/\mathfrak m^2$ (où $\mathfrak m=\mathfrak nA$) une structure
de $k$-extension de $K$ par $V$ (\emph{distincte} en général de la structure
$k$-triviale définie par l'injection canonique $k\to K\to A$) telle que $h$
soit l'homomorphisme caractéristique de cette extension. Si
$f:k\to A/\mathfrak m^2$ est l'homomorphisme structural, le fait que $k$
soit séparable sur le corps premier permet (19.6.1) de factoriser $f$ en
$k\xrightarrow{g}A\to A/\mathfrak m^2$, et la $k$-algèbre $A$ ainsi définie
répond à la question (20.6.22.2).
\end{proof}

\begin{theorem}[22.2.6]
\label{0.22.2.6-fr}
Soient $k$ un corps, $K$ une extension de $k$. Pour qu'il existe une
$k$-algèbre locale noethérienne $A$, d'idéal maximal $\mathfrak m$, telle que
$A/\mathfrak m=K$ et que $A$ soit formellement lisse (pour sa topologie
$\mathfrak m$-préadique), il faut et il suffit que $\Upsilon_{K/k}$ soit un
$K$-espace vectoriel de rang fini. La $k$-algèbre $A$ est alors déterminée
(à un $k$-isomorphisme près donnant par passage aux quotients l'identité
$K\to K$) par sa dimension, qui est soumise à la seule condition d'être
$\geqslant\operatorname{rg}_K\Upsilon_{K/k}$.
\end{theorem}

\begin{proof}
Comme $\mathfrak m/\mathfrak m^2$ est de rang fini sur $K$, la nécessité de
la condition résulte de (22.2.2, (ii))~; inversement, (22.2.5) montre que la
condition est suffisante et que l'on peut prendre pour
$\mathfrak m/\mathfrak m^2$ un $K$-espace vectoriel de rang quelconque
$\geqslant\operatorname{rg}_K(\Upsilon_{K/k})$.
\end{proof}

\oldpage[0\textsubscript{IV-1}]{189}

En particulier, l'homomorphisme identique de $\Upsilon_{K/k}$ associe
\emph{canoniquement} à $K$ une $k$-algèbre noethérienne complète,
formellement lisse, pour laquelle $\Upsilon_{K/k}$ est isomorphe à
$\mathfrak m/\mathfrak m^2$.

\begin{remarks}[22.2.7]
\label{0.22.2.7-fr}
\begin{enumerate}
  \item[(i)] Soit $\Lambda$ un anneau local noethérien, de corps résiduel
  $k$~; il résulte de (19.7.1) et (19.7.2) que la détermination des
  $\Lambda$-algèbres $A$ qui sont des anneaux locaux noethériens complets et
  qui sont formellement lisses est \emph{équivalente} au même problème où
  $\Lambda$ est remplacé par $k$, et est donc en principe entièrement
  résolue par (22.2.5), qui ramène la question à la structure des \emph{corps
  résiduels} des $k$-algèbres cherchées, en tant qu'extensions de $k$.

  \item[(ii)] Le fait que $\Upsilon_{K/k}$ soit de rang fini sur $K$
  \emph{n'entraîne pas} que $K$ soit «~de multiplicité radicielle finie~» sur
  $k$ (cf. (19.6.6)). Par exemple, soient $k_0$ un corps parfait de
  caractéristique $p>0$, $k=k_0(T)$ le corps des fractions rationnelles à une
  indéterminée sur $k_0$ et $K=k^{p^{-\infty}}$ la plus petite extension
  parfaite de $k$. On a alors $\Omega^1_K=\Omega^1_{K/k}=0$ (21.4.4), donc
  par définition (20.6.1.1)
  $\Upsilon_{K/k}=\Omega^1_k\otimes_kK$, et comme $k^p=k_0(T^p)$, $T$ est
  une $p$-base de $k$ sur $k^p$, donc (21.4.5), $\Omega^1_k$ est de rang $1$
  sur $k$, et par suite $\Upsilon_{K/k}$ de rang $1$ sur $K$. Mais pour tout
  $n$ le corps résiduel de l'anneau local $K\otimes_kk^{p^{-n}}$ est
  isomorphe à $K$ donc n'est pas séparable sur $k^{p^{-n}}$.
\end{enumerate}
\end{remarks}

Le théorème (22.2.2) se généralise comme suit~:

\begin{proposition}[22.2.8]
\label{0.22.2.8-fr}
Sous les conditions de (22.2.1), supposons que la caractéristique $p$ de $k$
soit $>0$~; soit $(k_\alpha)$ une famille filtrante décroissante de
sous-corps de $k$, telle que
$\bigcap_\alpha k_\alpha(k^p)=k^p$. Alors la condition (ii) de (22.2.2) peut
se remplacer par l'une quelconque des suivantes~:
\begin{enumerate}
  \item[a)] Il existe $\alpha_0$ tel que l'homomorphisme canonique
  \[
    \Omega^1_{k/k_\alpha}\otimes_kK
    \longrightarrow\Omega^1_{A/k_\alpha}\otimes_AK
  \]
  soit injectif pour tout $\alpha\geqslant\alpha_0$.
  \item[b)] Il existe $\alpha_0$ tel que $A$ soit une $k$-algèbre
  formellement lisse (pour la topologie $\mathfrak m$-préadique) relativement
  à $k_\alpha$, pour tout $\alpha\geqslant\alpha_0$.
  \item[c)] Pour tout $\alpha$ et tout entier $n\geqslant0$,
  l'homomorphisme canonique
  \[
    \Omega^1_{k/k_\alpha}\otimes_k(A/\mathfrak m^{n+1})
    \longrightarrow
    \Omega^1_{A/k_\alpha}\otimes_A(A/\mathfrak m^{n+1})
  \]
  est inversible à gauche.
\end{enumerate}
\end{proposition}

\begin{proof}
En effet, si $A$ est une $k$-algèbre formellement lisse, elle l'est \emph{a
fortiori} relativement à tout sous-corps de $k$, ce qui entraîne $c)$ en
vertu de (20.7.2) et (19.1.7)~; il est clair que $c)$ implique $b)$, en vertu
de (20.7.2), et que $b)$ implique $a)$. Enfin, compte tenu de (22.2.3), il
reste à prouver que $a)$ entraîne que l'homomorphisme canonique
\[
  u:\Omega^1_k\otimes_kK\longrightarrow\Omega^1_A\otimes_AK
\]
est injectif. Or, on a pour tout $\alpha$ un diagramme commutatif
\[
\xymatrix@C=5em@R=3em{
  \Omega^1_k\otimes_kK \ar[r]^{u} \ar[d]_{v_\alpha} &
  \Omega^1_A\otimes_AK \ar[d] \\
  \Omega^1_{k/k_\alpha}\otimes_kK \ar[r]_{u_\alpha} &
  \Omega^1_{A/k_\alpha}\otimes_AK
}
\]

\oldpage[0\textsubscript{IV-1}]{190}

et l'hypothèse que $u_\alpha$ est injectif pour
$\alpha\geqslant\alpha_0$ implique que $\operatorname{Ker}(u)$ est contenu
dans l'\emph{intersection} des $\operatorname{Ker}(v_\alpha)$. Mais on a vu
(21.8.3) que cette intersection est nulle en vertu de l'hypothèse
$\bigcap_\alpha k_\alpha(k^p)=k^p$, ce qui achève la démonstration.
\end{proof}

\begin{proposition}[22.2.9]
\label{0.22.2.9-fr}
Sous les hypothèses de (22.2.1), les conditions suivantes sont
équivalentes~:
\begin{enumerate}
  \item[a)] $A/\mathfrak m^2$ est une $k$-extension $k$-triviale de $K$ par
  $\mathfrak m/\mathfrak m^2$.
  \item[b)] $\chi_{A/k}=0$.
  \item[c)] L'homomorphisme canonique (20.5.11.2)
  \[
    \delta_{K/A/k}:\mathfrak m/\mathfrak m^2
    \longrightarrow\Omega^1_{A/k}\otimes_AK
  \]
  est injectif (en d'autres termes, la suite
  \begin{equation}
  \tag{22.2.9.1}
  \label{0.22.2.9.1-fr}
    0\longrightarrow\mathfrak m/\mathfrak m^2
    \longrightarrow\Omega^1_{A/k}\otimes_AK
    \longrightarrow\Omega^1_{K/k}\longrightarrow0
  \end{equation}
  est exacte).
\end{enumerate}
\end{proposition}

\begin{proof}
L'équivalence de $b)$ et $c)$ résulte de la suite exacte (20.6.22.1)~;
l'équivalence de $c)$ et $a)$ résulte de (20.5.12), la suite (22.2.9.1)
étant scindée si elle est exacte, puisqu'il s'agit d'espaces vectoriels.

On notera que si $K$ est \emph{séparable sur $k$}, les conditions
équivalentes de (22.2.9) sont satisfaites, en vertu de la définition de
$\chi_{A/k}$ (22.2.1) et de (20.6.3).
\end{proof}

\begin{corollary}[22.2.10]
\label{0.22.2.10-fr}
Sous les hypothèses de (22.2.1), considérons un sous-corps $k_0$ de $k$. Les
conditions suivantes sont équivalentes~:
\begin{enumerate}
  \item[a)] $A/\mathfrak m^2$ est une $k_0$-extension $k_0$-triviale de $K$
  par $\mathfrak m/\mathfrak m^2$.
  \item[b)] L'homomorphisme composé
  \[
    \Upsilon_{K/k_0}\longrightarrow\Upsilon_{K/k}
    \xrightarrow{\chi_{A/k}}\mathfrak m/\mathfrak m^2
  \]
  est nul.
  \item[c)] L'homomorphisme caractéristique
  $\chi_{A/k}:\Upsilon_{K/k}\to\mathfrak m/\mathfrak m^2$ se factorise en
  \[
    \Upsilon_{K/k}\longrightarrow\Upsilon_{K/k/k_0}
    \xrightarrow{\chi'}\mathfrak m/\mathfrak m^2.
  \]
\end{enumerate}
En outre, lorsque ces conditions sont remplies, l'homomorphisme $\chi'$ est
déterminé de façon unique et coïncide avec l'homomorphisme caractéristique
de la $k_0$-extension $k_0$-triviale $A/\mathfrak m^2$ de $K$ par
$\mathfrak m/\mathfrak m^2$.
\end{corollary}

\begin{proof}
Comme le composé dans $b)$ n'est autre que $\chi_{A/k_0}$, l'équivalence de
$a)$ et $b)$ résulte de (22.2.9)~; d'autre part, $b)$ et $c)$ sont
équivalentes en vertu de la suite exacte (20.6.18.1)
\[
  \Upsilon_{K/k_0}\longrightarrow\Upsilon_{K/k}
  \longrightarrow\Upsilon_{K/k/k_0}\longrightarrow0,
\]
et comme l'homomorphisme canonique
$\Upsilon_{K/k}\to\Upsilon_{K/k/k_0}$ est surjectif, cela implique l'unicité
de $\chi'$~; le fait que $\chi'$ est l'homomorphisme caractéristique de la
$k_0$-extension $A/\mathfrak m^2$ résulte de (20.6.22.2).
\end{proof}

\begin{corollary}[22.2.11]
\label{0.22.2.11-fr}
Sous les hypothèses de (22.2.1), soit $p$ l'exposant caractéristique de $k$,
et soit $(k_\alpha)$ une famille filtrante décroissante de sous-corps de $k$
telle que $\bigcap_\alpha k_\alpha(k^p)=k^p$. Si $\Upsilon_{K/k}$ est de rang
fini sur $K$, il existe un indice $\alpha$ tel que $A/\mathfrak m^2$ soit une
$k_\alpha$-extension $k_\alpha$-triviale de $K$ par
$\mathfrak m/\mathfrak m^2$.
\end{corollary}

\begin{proof}
On sait en effet (21.8.3) qu'il existe alors un indice $\alpha$ tel que
l'homomorphisme canonique
$\Upsilon_{K/k}\to\Upsilon_{K/k/k_\alpha}$ soit bijectif, donc la condition
$b)$ de (22.2.10) (où $k_\alpha$ remplace $k_0$) est satisfaite.
\end{proof}
\oldpage[0\textsubscript{IV-1}]{191}

\subsection{Application aux relations entre certains anneaux locaux et leurs complétés}
\label{subsection:0.22.3-fr}

\begin{lemma}[22.3.1]
\label{0.22.3.1-fr}
Soient $A_0$ un anneau semi-local noethérien, $A$ une $A_0$-algèbre finie
(qui est donc un anneau semi-local noethérien, cf. Bourbaki,
\emph{Alg. comm.}, chap.~IV, §~2, n\textsuperscript{o}~5, cor.~3 de la
prop.~9). Si $\widehat A_0$ et $\widehat A$ sont les complétés de $A_0$ et
$A$ pour leurs topologies préadiques respectives, alors, lorsqu'on munit
$A_0$, $A$ et $\widehat A$ des topologies discrètes, $\widehat A$ est une
$A$-algèbre formellement lisse relativement à $A_0$ (19.9.1).
\end{lemma}

\begin{proof}
On sait en effet que
$\widehat A=A\otimes_{A_0}\widehat A_0$ (Bourbaki,
\emph{Alg. comm.}, chap.~IV, §~2, n\textsuperscript{o}~5, cor.~3 de la
prop.~9 et chap.~III, §~3, n\textsuperscript{o}~4, th.~1), donc la proposition
résulte de (19.9.3).
\end{proof}

\begin{proposition}[22.3.2]
\label{0.22.3.2-fr}
Soient $A$ un anneau local noethérien régulier, $K$ son corps des fractions,
$p$ la caractéristique de $K$. On suppose vérifiée l'une des hypothèses
suivantes~:
\begin{enumerate}
  \item[(i)] $p=0$.
  \item[(ii)] $p>0$ et il existe une famille filtrante décroissante
  $(A_\alpha)$ de sous-anneaux noethériens de $A$, tels que $A$ soit une
  $A_\alpha$-algèbre finie pour tout $\alpha$, que
  $A^{p^{h(\alpha)}}\subset A_\alpha$ pour un entier $h(\alpha)$ convenable,
  et que, si $K_\alpha$ désigne le corps des fractions de $A_\alpha$, on ait
  \[
    \bigcap_\alpha K_\alpha(K^p)=K^p.
  \]
\end{enumerate}
Soient alors $B$ une $A$-algèbre finie intègre et contenant $A$,
$\mathfrak n$ un idéal premier de $B$, $C$ l'anneau local $B_\mathfrak n$,
$\mathfrak q$ un idéal premier du complété $\widehat C$ tel que
$\mathfrak q\cap C=0$, de sorte que l'anneau local
$\widehat C_\mathfrak q$ est une algèbre sur le corps des fractions $L$ de
$B$. Alors $\widehat C_\mathfrak q$ est une $L$-algèbre formellement lisse
pour sa topologie $\mathfrak q$-adique, et par suite un anneau
géométriquement régulier sur $L$ (19.6.5).
\end{proposition}

\begin{proof}
Distinguons deux cas.

I) Supposons que $\mathfrak n$ contienne l'idéal maximal $\mathfrak m$ de
$A$, et par suite $\mathfrak n\cap A=\mathfrak m$. Alors $\mathfrak n$ est
maximal (Bourbaki, \emph{Alg. comm.}, chap.~V, §~2,
n\textsuperscript{o}~1, prop.~1) et si $\mathfrak r$ est le radical de
l'anneau semi-local $B$, $\widehat C$ est le séparé complété de $B$ pour la
topologie $\mathfrak n$-préadique, et un des composants de l'anneau
semi-local $\widehat B$, complété de $B$ pour la topologie
$\mathfrak r$-préadique (Bourbaki, \emph{Alg. comm.}, chap.~III, §~3,
n\textsuperscript{o}~4, prop.~8)~; on peut donc écrire
$\widehat C_\mathfrak q=\widehat B_{\mathfrak q'}$, où $\mathfrak q'$ est
l'idéal premier de $\widehat B$ image réciproque de $\mathfrak q$. Notons
maintenant que $B$ est un sous-anneau de $\widehat B$ et l'hypothèse
$\mathfrak q\cap C=0$ entraîne $\mathfrak q'\cap B=0$, donc
$\widehat B_{\mathfrak q'}$ est aussi localisé de l'anneau
$\widehat B\otimes_B L$ en un de ses idéaux premiers $\mathfrak q''$, dont
$\mathfrak q'$ est l'image réciproque dans $\widehat B$. En outre on a
$\widehat B=\widehat A\otimes_A B$ (Bourbaki, \emph{Alg. comm.}, chap.~IV,
§~2, n\textsuperscript{o}~5, cor. de la prop.~9 et chap.~III, §~3,
n\textsuperscript{o}~4, th.~1), donc
$\widehat B\otimes_B L=\widehat A\otimes_A L$~; enfin, si $\mathfrak p$ est
l'idéal premier de $\widehat A$, image réciproque de $\mathfrak q''$, le
fait que $A\to B$ est injectif et la commutativité du diagramme
\[
  \xymatrix{
    B \ar[r] & \widehat B\\
    A \ar[u] \ar[r] & \widehat A \ar[u]
  }
\]
entraînent que $\mathfrak p\cap A=0$, donc
$(\widehat A\otimes_A L)_{\mathfrak q''}$ est aussi localisé de l'anneau
\[
  \widehat A_\mathfrak p\otimes_A L
  =(\widehat A_\mathfrak p\otimes_A K)\otimes_K L
\]
en un de ses idéaux premiers.

\oldpage[0\textsubscript{IV-1}]{192}

Pour montrer que cet anneau local est une $L$-algèbre formellement lisse,
il suffit donc de prouver que $\widehat A_\mathfrak p\otimes_A K$ est une
$K$-algèbre formellement lisse (19.3.5, (iii) et (iv))~; d'ailleurs
$\widehat A_\mathfrak p\otimes_A K=\widehat A_\mathfrak p$ puisque
$\mathfrak p\cap A=0$. On va appliquer les critères (22.2.2) et (22.2.8).
En premier lieu, comme $A$ est régulier, il en est de même de $\widehat A$
(17.1.5), donc de $\widehat A_\mathfrak p$ (17.3.2), et la condition (i) de
(22.2.2) est vérifiée (17.1.1). Il suffit par suite (lorsque $p>0$) de
vérifier la condition b) de (22.2.8). Or, (22.3.1) montre que, pour tout
$\alpha$, $\widehat A$ est une $A$-algèbre formellement lisse relativement
à $A_\alpha$ (pour les topologies discrètes)~; par suite
\[
  \widehat A\otimes_{A_\alpha}K_\alpha
  =\widehat A\otimes_A(A\otimes_{A_\alpha}K_\alpha)
\]
est une $(A\otimes_{A_\alpha}K_\alpha)$-algèbre formellement lisse
relativement à $K_\alpha$, pour les topologies discrètes (19.9.2, (iii)).
Mais comme $A$ est une $A_\alpha$-algèbre finie, on a
$A\otimes_{A_\alpha}K_\alpha=K$~; comme $\widehat A_\mathfrak p$ est un
anneau de fractions de $\widehat A\otimes_A K$, on en conclut (19.9.2,
(iv)) que $\widehat A_\mathfrak p$ est une $K$-algèbre formellement lisse
relativement à $K_\alpha$, pour la topologie discrète, et \emph{a fortiori}
pour sa topologie $\mathfrak p$-adique ((19.3.8) et (19.9.5)). Mais cela est
précisément la condition b) de (22.2.8), donc la proposition est démontrée
dans ce cas pour $p>0$. Lorsque $p=0$, le corps résiduel de
$\widehat A_\mathfrak p$ est séparable sur $K$, et comme
$\widehat A_\mathfrak p$ est un anneau régulier, c'est une $K$-algèbre
formellement lisse en vertu de (19.6.4).

II) Cas général. Soit $\mathfrak s$ l'idéal premier $\mathfrak n\cap A$ de
$A$, et posons $S=A-\mathfrak s$, de sorte que $S^{-1}A=A_\mathfrak s$, et
l'on a $C=(S^{-1}B)_{S^{-1}\mathfrak n}$, où cette fois
$S^{-1}\mathfrak n$ contient l'idéal maximal
$\mathfrak sA_\mathfrak s$ de $A_\mathfrak s$. Comme $A_\mathfrak s$ est
régulier (17.3.2) et $S^{-1}B$ une $A_\mathfrak s$-algèbre finie intègre
contenant $A_\mathfrak s$, tout revient à vérifier la condition (ii) pour
$A_\mathfrak s$ lorsque $p>0$~: or, posons
$\mathfrak s_\alpha=\mathfrak s\cap A_\alpha$ et considérons l'anneau local
$(A_\alpha)_{\mathfrak s_\alpha}$. Pour tout $t\notin\mathfrak s$, on a par
hypothèse, en posant $q=p^{h(\alpha)}$, $t^q\in A_\alpha$ et
$t^q\notin\mathfrak s$, donc $t^q\notin\mathfrak s_\alpha$~; si l'on pose
$S_\alpha=A_\alpha-\mathfrak s_\alpha$, on voit donc que l'on a
$A_\mathfrak s=S_\alpha^{-1}A$, et l'hypothèse entraîne que
$A_\mathfrak s$ est une $(A_\alpha)_{\mathfrak s_\alpha}$-algèbre finie~;
en outre, $K_\alpha$ est aussi le corps des fractions de
$(A_\alpha)_{\mathfrak s_\alpha}$, ce qui achève la démonstration.
\end{proof}

\begin{theorem}[22.3.3]
\label{0.22.3.3-fr}
Soient $A$ un anneau local noethérien complet, $\mathfrak p$ un idéal
premier de $A$, $B$ l'anneau localisé $A_\mathfrak p$, $\mathfrak q$ un
idéal premier du complété $\widehat B$, $\mathfrak r=\mathfrak q\cap B$.
Alors l'anneau localisé $\widehat B_\mathfrak q$ de $\widehat B$ est une
$B_\mathfrak r$-algèbre formellement lisse pour les topologies préadiques.
\end{theorem}

\begin{proof}
L'idéal premier $\mathfrak r$ de $B$ est de la forme
$\mathfrak n_\mathfrak p$, où $\mathfrak n\subset\mathfrak p$ est un idéal
premier de $A$, et l'on a $B_\mathfrak r=A_\mathfrak n$~; soit $k$ le corps
résiduel de $B_\mathfrak r$, qui est aussi le corps des fractions de l'anneau
local intègre et complet $A'=A/\mathfrak n$. Comme $\widehat B$ est un
$B$-module plat $(\mathbf 0_\mathrm I,~7.3.5)$,
$\widehat B_\mathfrak q$ est un $B_\mathfrak r$-module plat
$(\mathbf 0_\mathrm I,~6.3.2)$, et pour prouver que
$\widehat B_\mathfrak q$ est une $B_\mathfrak r$-algèbre formellement
lisse, il suffit de prouver que
\[
  \widehat B_\mathfrak q\otimes_{B_\mathfrak r}k
  =\widehat B_\mathfrak q/\mathfrak r\widehat B_\mathfrak q
\]
est une $k$-algèbre formellement lisse (19.7.1). Or
$\widehat B_\mathfrak q/\mathfrak r\widehat B_\mathfrak q
=(\widehat B/\mathfrak r\widehat B)_{\mathfrak q'}$, où $\mathfrak q'$ est
l'image canonique de $\mathfrak q$ dans
$\widehat B/\mathfrak r\widehat B$~; on a par définition
$\mathfrak q'\cap(B/\mathfrak r)=0$ et
$B/\mathfrak r=A_{\mathfrak p'}$, où
$\mathfrak p'=\mathfrak p/\mathfrak n$~; en outre
$\widehat B/\mathfrak r\widehat B$ est le complété de $B/\mathfrak r$.

On est donc ramené à démontrer le corollaire suivant~:
\end{proof}

\begin{corollary}[22.3.4]
\label{0.22.3.4-fr}
Soient $A$ un anneau local noethérien complet et intègre, $K$ son corps des
fractions, $\mathfrak p$ un idéal premier de $A$, $B$ l'anneau localisé
$A_\mathfrak p$. Alors, pour tout idéal premier $\mathfrak q$

\oldpage[0\textsubscript{IV-1}]{193}

du complété $\widehat B$, tel que $\mathfrak q\cap B=0$, l'anneau local
$\widehat B_\mathfrak q$ est une $K$-algèbre formellement lisse pour sa
topologie $\mathfrak q$-adique, donc un anneau géométriquement régulier sur
$K$.
\end{corollary}

\begin{proof}
En effet, il résulte de (19.8.9) qu'il existe un sous-anneau $A_0$ de $A$
qui est un anneau de séries formelles sur un corps de caractéristique $p>0$
ou sur un anneau de Cohen (on rappelle que les anneaux de Cohen ont un corps
des fractions de caractéristique $0$), et est tel que $A$ soit une
$A_0$-algèbre finie. Or, on sait que l'anneau $A_0$ est régulier (17.3.8) et
vérifie l'une des hypothèses (i), (ii) de (22.3.2), en vertu de (21.8.8)~;
il suffit donc d'appliquer (22.3.2), en remplaçant $B$ par $A$ et $A$ par
$A_0$.
\end{proof}

\subsection{Résultats préliminaires sur les extensions finies d'anneaux locaux dont l'idéal maximal est de carré nul}
\label{subsection:0.22.4-fr}
\label{0.22.4-fr}

\begin{proposition}[22.4.1]
\label{0.22.4.1-fr}
Soient $A$ un anneau local dont l'idéal maximal $\mathfrak m$ est de carré
nul, $K=A/\mathfrak m$ son corps résiduel~; désignons par $V$ l'idéal
$\mathfrak m$ considéré comme $K$-espace vectoriel. Soient $T_i$
$(1\leqslant i\leqslant r)$ des indéterminées, et pour chaque $i$, soit
$F_i$ un polynôme unitaire de $A[T_i]$ de degré $d_i$~; désignons par $A'$ le
quotient de l'anneau de polynômes $A[T_1,\ldots,T_r]$ par l'idéal
$\mathfrak b$ engendré par les $r$ polynômes $F_i$. Supposons que $A'$ soit
un anneau local, et soient $\mathfrak m'$ son idéal maximal,
$K'=A'/\mathfrak m'$ son corps résiduel. Alors~:
\begin{enumerate}
  \item[(i)] Si l'on pose $d=\prod_{i=1}^r d_i$, $A'$ est un $A$-module
  libre de rang $d$, et l'on a $[K':K]\leqslant d$.
  \item[(ii)] Pour que $[K':K]=d$, il faut et il suffit que
  $A'\otimes_AK=A'/\mathfrak mA'$ soit un corps (nécessairement isomorphe à
  $K'$), autrement dit, que $\mathfrak m'=\mathfrak mA'$. On a dans ce cas
  $\mathfrak m'^2=0$, et si $V'$ désigne l'idéal $\mathfrak m'$ considéré
  comme $K'$-espace vectoriel, l'homomorphisme canonique
  $V\otimes_KK'\to V'$ est bijectif (et par suite
  $\operatorname{rg}_{K'}V'=\operatorname{rg}_K V$).
\end{enumerate}
\end{proposition}

\begin{proof}
Sans supposer $A'$ local, il est évident, par division euclidienne et
récurrence sur $r$, que si $t_i$ est la classe de $T_i$ mod.~$\mathfrak b$,
les monômes
$M_\nu=t_1^{\nu_1}t_2^{\nu_2}\cdots t_r^{\nu_r}$, où
$0\leqslant\nu_i\leqslant d_i-1$, forment une base du $A$-module $A'$. Si
$A'$ est supposé local, $K'$ est un quotient de
$A'/\mathfrak mA'=A'\otimes_AK$, donc
$[K':K]\leqslant[A'\otimes_AK:K]=d$, et il ne peut y avoir égalité que si
$\mathfrak m'=\mathfrak mA'$, ce qui démontre (i) et la première assertion
de (ii). D'autre part, il est clair que lorsque
$\mathfrak m'=\mathfrak mA'$, on a $\mathfrak m'^2=0$, et
$V\otimes_KK'=\mathfrak m\otimes_AA'=\mathfrak mA'=V'$ puisque $A'$ est
un $A$-module libre.
\end{proof}

\begin{remark}[22.4.2]
\label{0.22.4.2-fr}
Lorsque $A$ est artinien (autrement dit $\operatorname{rg}_K(V)$ fini),
l'hypothèse que $A'$ est local entraîne toujours que
$\operatorname{rg}_{K'}(\mathfrak m'/\mathfrak m'^2)
\geqslant\operatorname{rg}_K(V)$, comme nous le verrons plus loin (22.5.2).
On notera que ces deux rangs peuvent être égaux sans que l'on ait
$\mathfrak m'=\mathfrak mA'$.
\end{remark}

\begin{proposition}[22.4.3]
\label{0.22.4.3-fr}
Les hypothèses sur $A$ et les notations étant celles de (22.4.1), supposons
que les $F_i$ soient de la forme
\begin{equation}
\label{0.22.4.3.1-fr}
  F_i(T_i)=T_i^{d_i}+\sum_{1\leqslant k\leqslant d_i}
  c_{ik}T_i^{d_i-k},
  \tag{22.4.3.1}
\end{equation}
où $d_i\geqslant2$ pour tout $i$, et $c_{ik}\in\mathfrak m$
(«~polynômes d'Eisenstein~»). Alors~:
\begin{enumerate}
  \item[(i)] $A'$ est un anneau local, et son corps résiduel
  $K'=A'/\mathfrak m'$ est isomorphe à $K$.

  \oldpage[0\textsubscript{IV-1}]{194}

  \item[(ii)] Le noyau de l'homomorphisme canonique
  $V\to V'=\mathfrak m'/\mathfrak m'^2$ est l'idéal $\mathfrak n$ de $A$
  engendré par les $\xi_i=c_{i,d_i}$ («~termes constants~» des $F_i$), et
  le conoyau de cet homomorphisme a pour base sur $K'$ les images des $T_i$~;
  en particulier, on a
  \begin{equation}
  \label{0.22.4.3.2-fr}
    \operatorname{rg}_{K'}(V')=\operatorname{rg}_K(V)+r-r',
    \tag{22.4.3.2}
  \end{equation}
  où $r'=\operatorname{rg}_K(\mathfrak n)$. Lorsque
  $\operatorname{rg}_K(V)$ est fini (autrement dit lorsque $A$ est
  artinien), pour que
  $\operatorname{rg}_{K'}(V')=\operatorname{rg}_K(V)$, il faut et il suffit
  que les $\xi_i$ soient linéairement indépendants sur $K$.
\end{enumerate}
\end{proposition}

\begin{proof}
(i) Il est clair que $A'\otimes_AK=A'/\mathfrak mA'$ est isomorphe à
$K[T_1,\ldots,T_r]/\mathfrak b'$, où $\mathfrak b'$ est l'idéal engendré
par les $r$ polynômes $T_i^{d_i}$~; c'est donc un anneau local de corps
résiduel $K$, dont l'idéal maximal est engendré par les classes des $T_i$
mod.~$\mathfrak b'$, d'où (i).

(ii) Soit $t_i$ $(1\leqslant i\leqslant r)$ la classe de $T_i$
mod.~$\mathfrak b$~; il résulte de (i) que l'idéal maximal $\mathfrak m'$ de
$A'$ est donné par
\[
  \mathfrak m'=\mathfrak mA'+\sum_{i=1}^r t_iA',
\]
d'où, compte tenu de $\mathfrak m^2=0$,
\[
  \mathfrak m'^2=\sum_{i,j}t_it_jA'+\sum_i t_i\mathfrak mA'.
\]
On en conclut que $A'/\mathfrak m'^2$ est isomorphe à l'anneau
$A[T_1,\ldots,T_r]/\mathfrak c$, où $\mathfrak c$ est l'idéal engendré par
les éléments
\[
  F_i\ (1\leqslant i\leqslant r),\qquad
  T_iT_j\ (1\leqslant i\leqslant r,\ 1\leqslant j\leqslant r),\qquad
  xT_i\ (1\leqslant i\leqslant r,\ x\in\mathfrak m).
\]
Comme $d_i\geqslant2$ et que les $T_i^2$ appartiennent à
$\mathfrak c$, les hypothèses faites entraînent que $\mathfrak c$ est aussi
engendré par les éléments
\[
  \xi_i\ (1\leqslant i\leqslant r),\qquad
  T_iT_j\ (1\leqslant i\leqslant r,\ 1\leqslant j\leqslant r),\qquad
  T_ix\ (1\leqslant i\leqslant r,\ x\in\mathfrak m).
\]
Soient
\[
  \mathfrak c_1=\sum_{i,j}T_iT_jA[T_1,\ldots,T_r],\qquad
  \mathfrak c_2=\mathfrak c_1+\sum_i\mathfrak mT_iA[T_1,\ldots,T_r],
\]
de sorte que
$\mathfrak c=\mathfrak c_2+\sum_i\xi_iA[T_1,\ldots,T_r]$. Il est clair que
$A[T_1,\ldots,T_r]/\mathfrak c_1$ est le $A$-module somme directe de $A$ et
des $At_i^{(1)}$, où $t_i^{(1)}$ est la classe de $T_i$
mod.~$\mathfrak c_1$. On en déduit que
$A[T_1,\ldots,T_r]/\mathfrak c_2$ est somme directe de $A$ et des
$Kt_i^{(2)}$, où $t_i^{(2)}$ est la classe de $T_i$
mod.~$\mathfrak c_2$. Enfin $A[T_1,\ldots,T_r]/\mathfrak c$ est somme
directe de $A/\mathfrak n$ et des $Kt_i'$, où $t_i'$ est la classe de $T_i$
mod.~$\mathfrak c$, et $\mathfrak n$ est l'idéal de $A$ engendré par les
$\xi_i$~; dans l'identification de $A'/\mathfrak m'^2$ à
$A[T_1,\ldots,T_r]/\mathfrak c$, l'image de $V$ s'identifie à
$\mathfrak m/\mathfrak n$, et ce qui précède prouve donc (ii).
\end{proof}

\begin{proposition}[22.4.4]
\label{0.22.4.4-fr}
Les hypothèses sur $A$ et les notations étant celles de (22.4.1), supposons
que $K$ soit de caractéristique $p>0$ (donc $p\cdot1\in\mathfrak m$ dans
$A$) et que les $F_i$ soient de la forme
\begin{equation}
\label{0.22.4.4.1-fr}
  F_i(T_i)=T_i^p+p\sum_{k=1}^{p-1}c_{ik}T_i^{p-k}-f_i
  \tag{22.4.4.1}
\end{equation}
avec $c_{ik}\in A$ et $f_i\in A$ pour tout $i$ et tout $k$ tel que
$1\leqslant k\leqslant p-1$~; en outre si $a_i$ est la classe de $f_i$
mod.~$\mathfrak m$, on suppose que la famille $(a_i)_{1\leqslant i\leqslant r}$
est $p$-libre sur $K^p$ («~absolument $p$-libre~»). Alors~:
\begin{enumerate}
  \item[(i)] $A'$ est un anneau local, d'idéal maximal $\mathfrak mA'$,
  dont le corps résiduel $K'$ est isomorphe à l'extension de $K$ obtenue par
  adjonction des $a_i^{1/p}$ $(1\leqslant i\leqslant r)$.

  \oldpage[0\textsubscript{IV-1}]{195}

  \item[(ii)] Les éléments $d_Af_i\otimes1_{K'}$ forment une base du noyau
  $N=\Upsilon^{K'}_{A'/A}$ de l'homomorphisme canonique
  \[
    \Omega^1_A\otimes_AK'\longrightarrow
    \Omega^1_{A'}\otimes_{A'}K'.
  \]
\end{enumerate}
\end{proposition}

\begin{proof}
(i) Cette fois $A'\otimes_AK=A'/\mathfrak mA'$ est isomorphe à
$K[T_1,\ldots,T_r]/\mathfrak b'$, où $\mathfrak b'$ est l'idéal engendré
par les polynômes $T_i^p-a_i$~; l'hypothèse faite sur les $a_i$ entraîne que
cet anneau quotient est un corps (21.1.8), d'où les assertions de (i).

Avant de démontrer (ii), nous établirons le

\begin{lemma}[22.4.4.2]
\label{0.22.4.4.2-fr}
Soient $\Lambda$ un anneau, $A$ une $\Lambda$-algèbre, $\mathfrak J$ un
idéal de $A$ tel que l'anneau $C=A/\mathfrak J$ soit de caractéristique
$p>0$ (ce qui équivaut à dire que $p\cdot1\in\mathfrak J$ dans $A$). Soit
$\pi$ l'image canonique de $p\cdot1$ dans
$\mathfrak J/\mathfrak J^2$. Si $C$ est une $(\Lambda/p\Lambda)$-algèbre
formellement lisse pour les topologies discrètes, on a une suite exacte
canonique
\begin{equation}
\label{0.22.4.4.3-fr}
  0\longrightarrow(\mathfrak J/\mathfrak J^2)/C\cdot\pi
  \longrightarrow\Omega^1_{A/\Lambda}\otimes_AC
  \longrightarrow\Omega^1_{C/\Lambda}\longrightarrow0.
  \tag{22.4.4.3}
\end{equation}
\end{lemma}

\begin{proof}[Démonstration du lemme]
Posons en effet $A'=A/pA$, $\mathfrak J'=\mathfrak J/pA$, de sorte que
$C=A'/\mathfrak J'$. Comme
$A'=A\otimes_\Lambda(\Lambda/p\Lambda)$, on a
$\Omega^1_{A/\Lambda}\otimes_AA'=\Omega^1_{A'/\Lambda'}$ à un
isomorphisme canonique près, $\Lambda'$ désignant l'anneau
$\Lambda/p\Lambda$ (20.5.5). Appliquant alors la suite exacte (20.5.14.1) à
la $\Lambda'$-algèbre formellement lisse $C=A'/\mathfrak J'$, on obtient la
suite exacte
\[
  0\longrightarrow\mathfrak J'/\mathfrak J'^2
  \longrightarrow\Omega^1_{A'/\Lambda'}\otimes_{A'}C
  \longrightarrow\Omega^1_{C/\Lambda'}\longrightarrow0
\]
et comme l'homomorphisme $\Lambda\to\Lambda'$ est surjectif, on a
$\Omega^1_{C/\Lambda'}=\Omega^1_{C/\Lambda}$ à un isomorphisme canonique
près (20.5.15)~; d'où la suite exacte (22.4.4.3) en vertu de ce qui précède.
\end{proof}

(ii) Comme on a $\mathfrak m'^2=0$, nous désignerons encore par $V'$ l'idéal
$\mathfrak m'$ considéré comme $K'$-espace vectoriel, et nous poserons
\begin{equation}
\label{0.22.4.4.4-fr}
  V_0=V/K\cdot p=\mathfrak m/(\mathfrak m^2+pA),\qquad
  V_0'=V'/K'\cdot p=\mathfrak m'/(\mathfrak m'^2+pA').
  \tag{22.4.4.4}
\end{equation}
Appliquant la suite exacte (22.4.4.3) au cas où $\Lambda=\mathbf Z$, à la
$\mathbf Z$-algèbre $A$ (resp. $A'$) et à son idéal $\mathfrak m$
(resp. $\mathfrak m'$), il vient, puisque $K$ (resp. $K'$) est séparable sur
le corps premier $\mathbf Z/p\mathbf Z$, donc une
$(\mathbf Z/p\mathbf Z)$-algèbre formellement lisse (19.6.1), les suites
exactes
\begin{equation}
\label{0.22.4.4.5-fr}
  0\longrightarrow V_0\longrightarrow\Omega^1_A\otimes_AK
  \longrightarrow\Omega^1_K\longrightarrow0,\qquad
  0\longrightarrow V_0'\longrightarrow\Omega^1_{A'}\otimes_{A'}K'
  \longrightarrow\Omega^1_{K'}\longrightarrow0.
  \tag{22.4.4.5}
\end{equation}
Considérons alors le diagramme (22.4.4.6) où les deux colonnes médianes sont
la seconde suite (22.4.4.5) et la première tensorisée par $K'$~; comme la
première suite exacte (22.4.4.5) est formée d'espaces vectoriels sur $K$,
les deux colonnes médianes de (22.4.4.6) sont exactes~; en outre, le
diagramme formé par ces colonnes et les flèches horizontales qui les
connectent est commutatif, en vertu de la démonstration de (22.4.4.2) et de
la commutativité du diagramme (20.5.11.3). La dernière ligne de (22.4.4.6)
est la suite exacte obtenue à partir de (20.6.1.1) en y remplaçant $A$, $B$,
$C$ par $\mathbf Z$, $K$ et $K'$~; la ligne médiane est obtenue par
tensorisation avec $K'$ de la suite exacte de $K$-modules déduite de
(20.6.1.1) en y remplaçant $A$, $B$, $C$ par $\mathbf Z$, $A$, $K$
respectivement. Le diagramme formé de ces deux lignes et des flèches
verticales qui les connectent est commutatif en vertu de la commutativité de
(20.6.4.3). Notons d'autre part qu'on est sous les conditions d'application
de (22.4.1, (ii)), donc la ligne supérieure du diagramme (22.4.4.6) est
exacte, les colonnes extrêmes du diagramme

\oldpage[0\textsubscript{IV-1}]{196}

\begin{equation}
\label{0.22.4.4.6-fr}
  \xymatrix{
    && 0 \ar[d] & 0 \ar[d] && \\
    & 0 \ar[r] & V_0\otimes_KK' \ar[r] \ar[d] & V_0' \ar[r] \ar[d] & 0 & \\
    0 \ar[r] & N \ar[r] \ar[d] & \Omega^1_A\otimes_AK' \ar[r] \ar[d] &
      \Omega^1_{A'}\otimes_{A'}K' \ar[r] \ar[d] &
      \Omega^1_{A'/A}\otimes_{A'}K' \ar[r] \ar[d] & 0 \\
    0 \ar[r] & \Upsilon_{K'/K} \ar[r] \ar[d] &
      \Omega^1_K\otimes_KK' \ar[r] \ar[d] & \Omega^1_{K'} \ar[r] \ar[d] &
      \Omega^1_{K'/K} \ar[r] \ar[d] & 0 \\
    & 0 & 0 & 0 & 0 &
  }
  \tag{22.4.4.6}
\end{equation}
sont donc formées respectivement des noyaux et des conoyaux des flèches
horizontales médianes, ce qui achève de prouver que le diagramme est
commutatif. En outre~:

\begin{lemma}[22.4.4.7]
\label{0.22.4.4.7-fr}
Les homomorphismes
\[
  N\longrightarrow\Upsilon_{K'/K}\qquad\text{et}\qquad
  \Omega^1_{A'/A}\otimes_{A'}K'\longrightarrow\Omega^1_{K'/K}
\]
du diagramme (22.4.4.6) sont bijectifs.
\end{lemma}

\begin{proof}
En effet, le diagramme du serpent (Bourbaki, \emph{Alg. comm.}, chap.~I,
§~1, n\textsuperscript{o}~4, prop.~2) appliqué aux deux colonnes médianes de
(22.4.4.6) donne une suite exacte
\[
  0\longrightarrow N\longrightarrow\Upsilon_{K'/K}\longrightarrow0
  \longrightarrow\Omega^1_{A'/A}\otimes_{A'}K'
  \longrightarrow\Omega^1_{K'/K}\longrightarrow0.
\]
Or, si $t_i$ est l'image canonique de $T_i$ dans $A'$, la relation
$F_i(t_i)=0$, jointe au fait que dans $A'$ on a
$p\cdot1\in\mathfrak mA'\subset\mathfrak m'$, montre aussitôt que l'on a
$d_{A'}f_i\in\mathfrak m'\Omega^1_{A'}$, donc
$d_{A'}f_i\otimes1_{K'}=0$~; cela prouve que les
$d_Af_i\otimes1_{K'}$ appartiennent à $N$. En outre, leurs images dans
$\Upsilon_{K'/K}$ sont les $d_Ka_i\otimes1_{K'}$, qui forment une base de
$\Upsilon_{K'/K}$ sur $K'$ (21.4.7)~; compte tenu de (22.4.4.7), cela achève
de démontrer (22.4.4).
\end{proof}
\end{proof}
\begin{env}[22.4.5]
\label{0.22.4.5-fr}
Considérons maintenant deux anneaux $A$, $B$ \emph{de caractéristique} $p>0$, et
deux homomorphismes
\[
  A\xrightarrow{i}B\xrightarrow{j}A
\]

\oldpage[0\textsubscript{IV-1}]{197}

vérifiant les conditions de (21.3.1). Soit en outre $\mathfrak m$ un idéal
\emph{de carré nul} dans $A$~; posons
\[
  K=A/\mathfrak m,\qquad
  B_{(K)}=B\otimes_AK=B/Bi(\mathfrak m)
\]
et soient
\[
  \varphi:A\longrightarrow K,\qquad
  \psi:B\longrightarrow B_{(K)}
\]
les homomorphismes canoniques.

Comme $\varphi$ est surjectif, $\Omega^1_{B_{(K)}/A}$ s'identifie
canoniquement à $\Omega^1_{B_{(K)}/K}$ (20.5.15). Notons que, puisque
$p\geqslant2$, pour tout $x\in\mathfrak m$, on a $j(i(x))=x^p=0$, autrement
dit $j(i(\mathfrak m))=0$~; par passage aux quotients on déduit donc de $j$ un
homomorphisme
\[
  j':B_{(K)}\longrightarrow A
\]
tel que, si $i'=\psi\circ i:A\longrightarrow B_{(K)}$, on ait
\[
  j'(i'(x))=x^p\qquad\text{et}\qquad i'(j'(y))=y^p
  \qquad\text{pour }x\in A\text{ et }y\in B_{(K)}.
\]
En posant, suivant les notations de (21.3.2),
\[
  \Theta_{B_{(K)}/A}
  =\Omega^1_{B_{(K)}/A}\otimes_{B_{(K)}}A_{[j']}
  =\Omega^1_{B_{(K)}/K}\otimes_{B_{(K)}}A_{[j']},
\]
on a donc (21.3.2.2) un homomorphisme canonique
\begin{equation}
\label{0.22.4.5.1-fr}
  \pi_{B_{(K)}/A}:\Theta_{B_{(K)}/A}\longrightarrow\Omega^1_A.
  \tag{22.4.5.1}
\end{equation}

D'autre part, on déduit de $i$ et $j$, par passage aux quotients (tenant compte
de ce que $j(i(\mathfrak m))=0$), des homomorphismes
\[
  K\xrightarrow{i_0}B_{(K)}\xrightarrow{j_0}K
\]
de sorte que l'on a le diagramme commutatif
\[
  \xymatrix{
    A \ar[r]^i \ar[d]_\varphi & B \ar[r]^j \ar[d]^\psi & A \ar[d]^\varphi \\
    K \ar[r]_{i_0} & B_{(K)} \ar[r]_{j_0} \ar[ur]_{j'} & K
  }
\]
Il est clair en outre que l'on a
\[
  j_0(i_0(s))=s^p\qquad\text{et}\qquad i_0(j_0(t))=t^p
  \qquad\text{pour }s\in K\text{ et }t\in B_{(K)}.
\]
On définit donc de même
\[
  \Theta_{B_{(K)}/K}
  =\Omega^1_{B_{(K)}/K}\otimes_{B_{(K)}}K_{[j_0]}
  =\Theta_{B_{(K)}/A}\otimes_AK,
\]
donc, en tensorisant (22.4.5.1) par $K$, on obtient un
$K$-\emph{homomorphisme canonique}
\begin{equation}
\label{0.22.4.5.2-fr}
  \pi'_{B_{(K)}/K}=\pi_{B_{(K)}/A}\otimes1_K:
  \Theta_{B_{(K)}/K}\longrightarrow\Omega^1_A\otimes_AK.
  \tag{22.4.5.2}
\end{equation}

\oldpage[0\textsubscript{IV-1}]{198}

Il résulte aussitôt des définitions que le diagramme
\[
  \xymatrix{
    \Theta_{B_{(K)}/K} \ar[r]^{\pi'_{B_{(K)}/K}} \ar[dr]_{\pi_{B_{(K)}/K}}
      & \Omega^1_A\otimes_AK \ar[d]^{\varphi_{K/A}} \\
    & \Omega^1_K
  }
\]
est commutatif, $\pi_{B_{(K)}/K}$ étant l'homomorphisme canonique
correspondant à $i_0$ et $j_0$ (21.3.2.2). Par \emph{restriction} à
$\Xi_{B_{(K)}/K}=\operatorname{Ker}(\pi_{B_{(K)}/K})$, l'homomorphisme
$\pi'_{B_{(K)}/K}$ définit donc un \emph{homomorphisme canonique}
\begin{equation}
\label{0.22.4.5.3-fr}
  \Xi_{B_{(K)}/K}\longrightarrow
  \operatorname{Ker}(\Omega^1_A\otimes_AK\longrightarrow\Omega^1_K)=\Upsilon_{K/A}.
  \tag{22.4.5.3}
\end{equation}
\end{env}

\begin{env}[22.4.6]
\label{0.22.4.6-fr}
Les hypothèses et notations étant celles de (22.4.5), supposons en outre que
$\mathfrak m$ soit un idéal \emph{maximal de carré nul}, donc $K$ un
\emph{corps}, et notons $V$ l'idéal $\mathfrak m$ considéré comme $K$-espace
vectoriel. Comme $K$ est une algèbre formellement lisse sur son corps premier
$\mathbf F_p$ (21.5.2) et que $A$ est une algèbre sur $\mathbf F_p$ (de sorte
que $\Omega^1_A=\Omega^1_{A/\mathbf F_p}$), on déduit de (20.5.14) que l'on a
une \emph{suite exacte}
\begin{equation}
\label{0.22.4.6.1-fr}
  0\longrightarrow V\longrightarrow\Omega^1_A\otimes_AK
  \longrightarrow\Omega^1_K\longrightarrow0.
  \tag{22.4.6.1}
\end{equation}
On a donc par (22.4.5.3) un homomorphisme noté
\begin{equation}
\label{0.22.4.6.2-fr}
  \chi'_{B/A}:\Xi_{B_{(K)}/K}\longrightarrow V
  \tag{22.4.6.2}
\end{equation}
que nous appellerons \emph{homomorphisme caractéristique} de la $A$-algèbre
$B$ (relative à $i$, $j$ et $\mathfrak m$). Il résulte des définitions que pour
que $\chi'_{B/A}$ soit \emph{injectif}, il faut et il suffit que
$\pi'_{B_{(K)}/K}$ soit injectif, le noyau de ce dernier étant évidemment
contenu dans $\Xi_{B_{(K)}/K}$.
\end{env}

\begin{proposition}[22.4.7]
\label{0.22.4.7-fr}
Soient $A$ un anneau local artinien dont l'idéal maximal $\mathfrak m$ est de
carré nul, $K=A/\mathfrak m$ son corps résiduel, $V$ l'idéal $\mathfrak m$
considéré comme $K$-espace vectoriel~; on suppose que $A$ est de
caractéristique $p>0$~; soient
\begin{equation}
\label{0.22.4.7.1-fr}
  F_i(T_i)=T_i^p-f_i\qquad\text{avec }f_i\in A\quad(1\leqslant i\leqslant r)
  \tag{22.4.7.1}
\end{equation}
et désignons par $B$ l'anneau quotient de $A[T_1,\ldots,T_r]$ par l'idéal
engendré par les $r$ polynômes $F_i$. Alors~:
\begin{enumerate}
  \item[\rm(i)] $B$ est un anneau local.
  \item[\rm(ii)] Si $\mathfrak m'$ est l'idéal maximal de $B$,
  $K'=B/\mathfrak m'$ son corps résiduel, $V'$ le $B$-module
  $\mathfrak m'/\mathfrak m'^2$ considéré comme $K'$-espace vectoriel, alors,
  pour que $\operatorname{rg}_K(V)=\operatorname{rg}_{K'}(V')$, il faut et il suffit que l'homomorphisme
  caractéristique de $K$-espaces vectoriels
  \begin{equation}
  \label{0.22.4.7.2-fr}
    \chi'_{B/A}:\Xi_{B_{(K)}/K}\longrightarrow V
    \tag{22.4.7.2}
  \end{equation}
  (cf. (22.4.6.2)) soit injectif.
\end{enumerate}
\end{proposition}

\begin{proof}
Désignons encore par $a_i$ $(1\leqslant i\leqslant r)$ la classe de $f_i$
mod.~$\mathfrak m$~; les $a_i$ ne sont plus nécessairement $p$-indépendants
sur $K^p$, mais on peut toujours supposer que la sous-famille

\oldpage[0\textsubscript{IV-1}]{199}

$(a_i)_{s+1\leqslant i\leqslant r}$ est $p$-libre sur $K^p$ et forme une
$p$-base sur $K^p$ du corps $K^p(a_1,\ldots,a_r)$ de sorte que l'on a
\[
  a_i\in K^p(a_{s+1},\ldots,a_r)\qquad\text{pour }1\leqslant i\leqslant s.
\]
Désignons par $A''$ l'anneau quotient de $A[T_{s+1},\ldots,T_r]$ par l'idéal
engendré par les $F_i$ d'indice $i\geqslant s+1$~; alors, par (22.4.4), $A''$
est un anneau local dont l'idéal maximal est $\mathfrak m''=\mathfrak mA''$
(donc de carré nul) et le corps résiduel
$K''=K(a_{s+1}^{1/p},\ldots,a_r^{1/p})$. On a donc
$a_i\in K''^p$ pour $1\leqslant i\leqslant s$ et par suite, dans $A''$,
l'élément $f_i$, pour $1\leqslant i\leqslant s$, peut s'écrire
\begin{equation}
\label{0.22.4.7.3-fr}
  f_i=g_i^p+h_i
  \tag{22.4.7.3}
\end{equation}
où $g_i\in A''$ et $h_i\in\mathfrak m''$ (comme $\mathfrak m''^2=0$ et
$p\mathfrak m''=0$, il est immédiat que $h_i$ est déterminé de façon
\emph{unique} par ces conditions). Remplaçant $T_i$ par $T_i+g_i$, on voit donc
que si l'on pose
\[
  G_i(T_i)=T_i^p+\sum_{k=1}^{p-1}\binom pk g_i^{p-k}T_i^k-h_i
  \qquad(1\leqslant i\leqslant s),
\]
$B$ est l'anneau quotient de $A''[T_1,\ldots,T_s]$ par l'idéal engendré par les
$G_i$~; or, tous les coefficients de $G_i$ sauf le coefficient dominant
appartiennent à $\mathfrak m''$, donc on est dans la situation de (22.4.3),
$B$ est un anneau local, et son corps résiduel $K'$ est \emph{isomorphe} à
$K''$, ce qui prouve (i) (qui d'ailleurs résulte directement de ce que
$B^p\subset A$, et par suite il n'y a qu'un seul idéal de $B$ au-dessus de
$\mathfrak m$).

Notons maintenant que pour que (22.4.7.2) soit injectif, il faut et il suffit
qu'il en soit ainsi de
\[
  \pi'_{B_{(K)}/K}:\Omega^1_{B_{(K)}/K}\otimes_{B_{(K)}}K
  \longrightarrow\Omega^1_A\otimes_AK
\]
(par (22.4.6)). Or $B_{(K)}=B/\mathfrak mB$ est l'algèbre quotient de
$C=K[T_1,\ldots,T_r]$ par l'idéal engendré par les polynômes
$\overline F_i=T_i^p-a_i$, et comme $d_{C/K}\overline F_i=0$ il résulte de
(20.5.13) que $\Omega^1_{B_{(K)}/K}$ est un $B_{(K)}$-module libre ayant pour
\emph{base} les $d_{B_{(K)}/K}t_i$, où $t_i$ est l'image canonique de $T_i$
$(1\leqslant i\leqslant r)$. Mais par définition (22.4.5) l'image canonique par
$j':B_{(K)}\longrightarrow A$ d'une classe mod.~$\mathfrak mB$ d'un élément
$x\in B$ est l'élément $x^p\in A$~; on en déduit aussitôt, puisque $t_i^p=f_i$,
que l'image de $d_{B_{(K)}/K}t_i\otimes1$ par $\pi'_{B_{(K)}/K}$ est
$d_Af_i\otimes1$ dans $\Omega^1_A\otimes_AK$~; la condition que
$\pi'_{B_{(K)}/K}$ soit injectif est donc équivalente au fait que les
$d_Af_i\otimes1_K$ sont \emph{linéairement indépendants sur $K$}, ou encore
que leurs images $d_Af_i\otimes1_{K'}$ sont \emph{linéairement indépendants
sur $K'$} dans $\Omega^1_A\otimes_AK'$.

Or, si l'on applique (22.4.4) à la $A$-algèbre $A''$, on voit que le noyau de
l'homomorphisme canonique
$\Omega^1_A\otimes_AK'\longrightarrow\Omega^1_{A''}\otimes_{A''}K'$ a pour
base les $d_Af_i\otimes1_{K'}$ pour $s+1\leqslant i\leqslant r$. La condition
précédente équivaut donc encore à dire que les $d_{A''}f_i\otimes1_{K'}$ sont
linéairement indépendants dans $\Omega^1_{A''}\otimes_{A''}K'$ pour
$1\leqslant i\leqslant s$. Or, on a $d_{A''}f_i=d_{A''}h_i$ par (22.4.7.3)~;
d'autre part, si $V''$ est le $K'$-espace vectoriel $\mathfrak m''$, on a
$\operatorname{rg}_{K'}(V'')=\operatorname{rg}_K(V)$ (22.4.1, (ii)), et la condition
$\operatorname{rg}_{K'}(V')=\operatorname{rg}_K(V)$ est donc équivalente à
$\operatorname{rg}_{K'}(V'')=\operatorname{rg}_{K'}(V')$. Mais il résulte de (22.4.3) que cette dernière
relation est équivalente au fait que les classes $h_i$

\oldpage[0\textsubscript{IV-1}]{200}

dans $V''$ sont \emph{linéairement indépendantes sur $K'$}. Or, dans
$\Omega^1_{A''}\otimes_{A''}K'$, les $d_{A''}h_i\otimes1_{K'}$ sont les images
canoniques (20.5.11.2) des $h_i$ par l'injection
$V''\longrightarrow\Omega^1_{A''}\otimes_{A''}K'$ (cf. (22.4.6.1)), et cela
achève de prouver (22.4.7).
\end{proof}

On a en outre prouvé~:

\begin{corollary}[22.4.7.4]
\label{0.22.4.7.4-fr}
Pour que $\operatorname{rg}_K(V)=\operatorname{rg}_{K'}(V')$, il faut et il suffit que les éléments
$d_Af_i\otimes1_K$ $(1\leqslant i\leqslant r)$ soient linéairement
indépendants dans le $K$-espace vectoriel $\Omega^1_A\otimes_AK$.
\end{corollary}

\begin{remark}[22.4.8]
\label{0.22.4.8-fr}
On peut définir les homomorphismes (22.4.5.2) et (22.4.6.2) dans des conditions
un peu différentes. Soient $A$ un anneau, $\mathfrak m$ un idéal \emph{de carré
nul} dans $A$, $K=A/\mathfrak m$, $p$ un nombre premier tel que
$p\cdot1\in\mathfrak m$ (autrement dit $K$ est de caractéristique $p$, mais non
nécessairement $A$). Soit d'autre part $B$ une $A$-algèbre qui est un $A$-module
\emph{fidèlement plat} (de sorte que $A\subset B$), et supposons que l'on ait
\begin{equation}
\label{0.22.4.8.1-fr}
  y^p\in A+pB\subset A+\mathfrak mB
  \tag{22.4.8.1}
\end{equation}
pour tout $y\in B$. Si $y^p=x+pz$ avec $x\in A$, $z\in B$, la \emph{classe}
mod.~$(A\cap pB)$ de $x$ est bien déterminée, et comme $A\cap pB=pA$ puisque
$B$ est un $A$-module fidèlement plat, on a ainsi défini une application
canonique $j:B\longrightarrow A/pA$, telle que, pour $x\in A$, $j(x)$ soit la
classe de $x^p$ mod.~$pA$. D'ailleurs, si $y'\in B$ et
$y'^p=x'+pz'$ avec $x'\in A$, $z'\in B$, il est immédiat que les éléments
$(y+y')^p-x-x'$ et $(yy')^p-xx'$ appartiennent à $pB$ en vertu du fait que les
coefficients binomiaux $\binom pi$ sont des multiples de $p$ pour
$1\leqslant i\leqslant p-1$. L'application $j$ est donc un
\emph{homomorphisme d'anneaux}. Par composition, on déduit de $j$ un
homomorphisme $j':B\xrightarrow{j}A/pA\longrightarrow A/\mathfrak m=K$ et par
suite un homomorphisme $j_0:B_{(K)}\longrightarrow K$ tel que $j'$ soit le
composé $B\longrightarrow B_{(K)}=B\otimes_AK\xrightarrow{j_0}K$~; en outre,
si $i_0:K\longrightarrow B_{(K)}$ est l'application canonique, $i_0$ et $j_0$
vérifient les conditions de (21.3.1) pour les anneaux \emph{de caractéristique
$p$}, $K$ et $B_{(K)}$. Cela étant, pour définir un $K$-homomorphisme
$\pi'_{B_{(K)}/K}:\Omega^1_{B_{(K)}/K}\otimes_{B_{(K)}}K_{[j_0]}
\longrightarrow\Omega^1_A\otimes_AK$, il suffit (20.4.8) de définir une
$K$-dérivation
\begin{equation}
\label{0.22.4.8.2-fr}
  D_0:B_{(K)}\longrightarrow\Omega^1_A\otimes_AK
  \tag{22.4.8.2}
\end{equation}
où le second membre est considéré comme un $B_{(K)}$-module au moyen de $j_0$.
Pour cela, il suffit de définir une $A$-dérivation
\begin{equation}
\label{0.22.4.8.3-fr}
  D:B\longrightarrow\Omega^1_A\otimes_AK
  \tag{22.4.8.3}
\end{equation}
qui \emph{s'annule dans} $\mathfrak mB$ (le second membre étant considéré comme
$B$-module au moyen de $j'$). Or, si l'on a la relation $y^p=x+pz$ avec
$y\in B$, $x\in A$, $z\in B$, l'élément $d_Ax\otimes1_K$ de
$\Omega^1_A\otimes_AK$ ne dépend que de la classe $j(y)$ mod.~$pA$ de $x$
(donc seulement de $y$), car pour tout $t\in A$ on a
$d_A(pt)=p\,d_A(t)$, donc $d_A(pt)\otimes1=0$ dans
$\Omega^1_A\otimes_AK=\Omega^1_A/\mathfrak m\Omega^1_A$, puisque
$p\cdot1\in\mathfrak m$. On peut donc prendre
\[
  D(y)=d_Ax\otimes1_K.
\]
Le fait que $D$ est une dérivation résulte aisément de ce que $j$ est un
homomorphisme d'anneaux~; en outre, pour $t\in A$, on a
$d_A(t^px)=t^p d_A(x)+pt^{p-1}x d_A(t)$ et l'on en tire que
$D(ty)=tD(y)$, donc $D$ est une $A$-dérivation~; enfin, si l'on a de plus
$t\in\mathfrak m$, on a

\oldpage[0\textsubscript{IV-1}]{201}

$D(ty)=tD(y)=0$ puisque $\Omega^1_A\otimes_AK$ est annulé par $\mathfrak m$.
On a ainsi défini l'homomorphisme (22.4.5.2) cherché, et il est immédiat de
vérifier que l'homomorphisme composé
\[
  \Theta_{B_{(K)}/K}\xrightarrow{\pi'_{B_{(K)}/K}}
  \Omega^1_A\otimes_AK\longrightarrow\Omega^1_K
\]
est encore l'homomorphisme canonique $\pi_{B_{(K)}/K}$ (relatif à $i_0$ et
$j_0$). On a donc aussi un homomorphisme canonique analogue à (22.4.5.3).

Si maintenant $K$ est un \emph{corps}, $A/pA$ est une $\mathbf F_p$-algèbre,
et l'on a la suite exacte (20.5.14)
\[
  0\longrightarrow V/pA\longrightarrow
  \Omega^1_{A/pA}\otimes_{A/pA}K\longrightarrow\Omega^1_K\longrightarrow0
\]
et d'autre part on a la suite exacte (20.5.12.1)
\[
  pA\longrightarrow\Omega^1_A\otimes_A(A/pA)
  \longrightarrow\Omega^1_{A/pA}\longrightarrow0
\]
qui montre, en tensorisant avec $K$, que les $K$-espaces vectoriels
$\Omega^1_A\otimes_AK$ et
$\Omega^1_{A/pA}\otimes_{A/pA}K$ sont isomorphes. L'analogue de (22.4.6.2) est
donc ici un homomorphisme
\begin{equation}
\label{0.22.4.8.4-fr}
  \Xi_{B_{(K)}/K}\longrightarrow V/pA=\mathfrak m/(\mathfrak m^2+pA).
  \tag{22.4.8.4}
\end{equation}

Ceci dit, \emph{le critère (22.4.7), où l'on remplace $V$ et $V'$ par $V/pA$
et $V'/pB$ respectivement, reste valable en supposant seulement que $K$ soit
de caractéristique $p>0$} (autrement dit, que $p\cdot1\in\mathfrak m$ dans
$A$), \emph{et que les $F_i(T_i)$ soient de la forme plus générale
(22.4.4.1)}~: en effet, $B$ est un $A$-module libre (donc fidèlement plat), et
il suffit de reprendre le raisonnement de (22.4.7), en y remplaçant (22.4.6.2)
par (22.4.8.4) et $V''$ par $V''/pA''$.
\end{remark}

\subsection{Algèbres géométriquement régulières et algèbres formellement lisses}
\label{subsection:0.22.5-fr}

\begin{proposition}[22.5.1]
\label{0.22.5.1-fr}
Soient $A$ un anneau local régulier, $A'$ un anneau local contenant $A$ et qui
est une $A$-algèbre finie, $\mathfrak m$ (resp. $\mathfrak m'$) l'idéal maximal
de $A$ (resp. $A'$), $K=A/\mathfrak m$ (resp. $K'=A'/\mathfrak m'$) son corps
résiduel. Pour que $A'$ soit un anneau régulier, il faut et il suffit que l'on
ait
\begin{equation}
\label{0.22.5.1.1-fr}
  \operatorname{rg}_K(\mathfrak m/\mathfrak m^2)
  =\operatorname{rg}_{K'}(\mathfrak m'/\mathfrak m'^2).
  \tag{22.5.1.1}
\end{equation}
\end{proposition}

\begin{proof}
En effet, le premier membre de (22.5.1.1) est $\dim(A)$ (17.1.1) et comme
$A\subset A'$ et que $A'$ est une $A$-algèbre finie,
$\dim(A')=\dim(A)$ (16.1.5)~; l'égalité (22.5.1.1) est donc nécessaire et
suffisante pour que $A'$ soit régulier (17.1.1).
\end{proof}

\begin{remark}[22.5.2]
\label{0.22.5.2-fr}
\begin{enumerate}
  \item[\rm(i)] Soient $A_1=A/\mathfrak m^2$,
  $A'_1=A'\otimes_AA_1=A'/\mathfrak m^2A'$~; comme
  $\mathfrak m^2A'\subset\mathfrak m'^2$,
  $\mathfrak m'_1=\mathfrak m'/\mathfrak m^2A'$ est l'idéal maximal de $A'_1$
  et $\mathfrak m'_1/\mathfrak m_1'^2$ est un $K'$-espace vectoriel isomorphe
  à $\mathfrak m'/\mathfrak m'^2$~; la condition de régularité de $A'$ ne
  dépend donc que de la structure de la $A_1$-algèbre $A'_1$, ce qui va nous
  permettre ci-dessous d'appliquer les résultats préliminaires de (22.4) sur
  les algèbres finies sur des anneaux artiniens.
  \item[\rm(ii)] Il résulte de la démonstration de (22.5.1) et de (16.2.6.1)
  que l'on a en tout cas
  \begin{equation}
  \label{0.22.5.2.1-fr}
    \operatorname{rg}_K(\mathfrak m/\mathfrak m^2)
    \leqslant\operatorname{rg}_{K'}(\mathfrak m'/\mathfrak m'^2).
    \tag{22.5.2.1}
  \end{equation}

  \oldpage[0\textsubscript{IV-1}]{202}

  \item[\rm(iii)] Supposons que $A$ et $A'$ vérifient les hypothèses générales
  de (22.4.1). On sait (19.8.10) qu'il existe un anneau local \emph{régulier}
  $B$, d'idéal maximal $\mathfrak n$, tel que $A$ soit isomorphe à
  $B/\mathfrak n^2$~; désignons par $G_i\in B[T_i]$ un polynôme unitaire dont
  l'image canonique dans $A[T_i]$ soit $F_i$ $(1\leqslant i\leqslant r)$~;
  alors, si $B'$ est l'anneau quotient de $B[T_1,\ldots,T_r]$ par l'idéal
  engendré par les $G_i$, il est clair que $B'$ est un $B$-module libre de rang
  $d$ et que $A'=B'\otimes_BA=B'/\mathfrak n^2B'$~; comme $A'$ est supposé être
  un anneau local, il en est de même de $B'$, et si $\mathfrak n'$ est l'idéal
  maximal de $B'$, $B'/\mathfrak n'$ est isomorphe à $K'$ et
  $\mathfrak n'/\mathfrak n'^2$ isomorphe à $\mathfrak m'/\mathfrak m'^2$ en
  tant qu'espace vectoriel sur $K'$~; l'inégalité (22.5.2.1) montre donc que
  l'on a, sous les hypothèses de (22.4.1),
  \begin{equation}
  \label{0.22.5.2.2-fr}
    \operatorname{rg}_K(V)\leqslant\operatorname{rg}_{K'}(\mathfrak m'/\mathfrak m'^2).
    \tag{22.5.2.2}
  \end{equation}
\end{enumerate}
\end{remark}

\begin{corollary}[22.5.3]
\label{0.22.5.3-fr}
Soient $A$ un anneau local régulier, $\mathfrak m$ son idéal maximal,
$T_i$ $(1\leqslant i\leqslant r)$ des indéterminées, et pour chaque $i$, soit
\begin{equation}
\label{0.22.5.3.1-fr}
  F_i(T_i)=T_i^{d_i}+\sum_{1\leqslant k\leqslant d_i}c_{ik}T_i^{d_i-k}
  \tag{22.5.3.1}
\end{equation}
un polynôme unitaire de $A[T_i]$, tel que $d_i\geqslant2$ pour tout $i$ et
$c_{ik}\in\mathfrak m$. Soit $A'$ le quotient de l'anneau de polynômes
$A[T_1,\ldots,T_r]$ par l'idéal engendré par les $r$ polynômes $F_i$. Alors~:
\begin{enumerate}
  \item[\rm(i)] $A'$ est un anneau local, et son corps résiduel $K'$ est
  isomorphe au corps résiduel $K$ de $A$.
  \item[\rm(ii)] Pour que $A'$ soit un anneau régulier, il faut et il suffit
  que les classes $\xi_i$ mod.~$\mathfrak m^2$ des éléments $c_{i,d_i}$
  $(1\leqslant i\leqslant r)$ soient linéairement indépendantes sur $K$.
\end{enumerate}
\end{corollary}

\begin{proof}
En effet, avec les notations de (22.5.2, (i)), si $\overline F_i$ désigne le
polynôme de $A_1[T_i]$ obtenu en appliquant aux coefficients de $F_i$
l'homomorphisme $A\longrightarrow A_1$, $A'_1$ se définit à partir de $A_1$ et
des polynômes $\overline F_i$ comme dans (22.4.1), et la conclusion résulte
donc de (22.5.1) et de (22.4.3).
\end{proof}

\begin{theorem}[22.5.4]
\label{0.22.5.4-fr}
Soient $A$ un anneau local régulier, $K$ son corps résiduel~; on suppose que
$K$ est de caractéristique $p>0$~; soient
\begin{equation}
\label{0.22.5.4.1-fr}
  F_i(T_i)=T_i^p+p\sum_{k=1}^{p-1}c_{ik}T_i^{p-k}-f_i
  \tag{22.5.4.1}
\end{equation}
avec $c_{ik}\in A$ et $f_i\in A$ $(1\leqslant i\leqslant r)$~; soit $B$
l'anneau quotient de $A[T_1,\ldots,T_r]$ par l'idéal engendré par les $r$
polynômes $F_i$. Alors $B$ est un anneau local, et les conditions suivantes
sont équivalentes~:
\begin{enumerate}
  \item[a)] $B$ est régulier~;
  \item[b)] les éléments $d_Af_i\otimes1_K$ $(1\leqslant i\leqslant r)$ sont
  linéairement indépendants dans le $K$-espace vectoriel
  $\Omega^1_A\otimes_AK$~;
  \item[c)] si $B'=B/\mathfrak m^2B$ et si l'on pose
  $\Xi_{B_{(K)}/K}=\Xi_{B'_{(K)}/K}$ (22.4.5), l'homomorphisme
  caractéristique $\Xi_{B_{(K)}/K}\longrightarrow\mathfrak m/\mathfrak m^2$
  (22.4.6.2) est injectif.
\end{enumerate}
\end{theorem}

\begin{proof}
En effet, $B'$ se définit encore à partir de $A_1=A/\mathfrak m^2$ et des
polynômes $\overline F_i$ obtenus en appliquant aux $F_i$ l'homomorphisme
$A\longrightarrow A_1$. Compte tenu de ce que $\Omega^1_A\otimes_AK$ est
canoniquement isomorphe à $\Omega^1_{A_1}\otimes_{A_1}K$ (20.5.12, (ii)), le
théorème résulte de (22.5.1) et de (22.4.7.4), en utilisant la remarque
(22.4.8) lorsque $A_1$ n'est pas de caractéristique $p$.
\end{proof}

\oldpage[0\textsubscript{IV-1}]{203}

\begin{env}[22.5.5]
\label{0.22.5.5-fr}
Soient $k$ un corps de caractéristique $p>0$, $A$ un anneau local, $k'$ une
extension \emph{radicielle finie} de $k$ telle que $k'^p\subset k$,
$\mathfrak m$ l'idéal maximal de $A$, $K=A/\mathfrak m$ son corps résiduel~;
rappelons que $A'=A\otimes_kk'$ est un anneau local dont le corps résiduel est
le même que celui de l'anneau local $K\otimes_kk'$ (Bourbaki,
\emph{Alg. comm.}, chap.~V, §~2, n\textsuperscript{o}~3, lemme~4). Notons que
l'on a $A'_{(K)}=A'\otimes_AK=K\otimes_kk'$, et par suite (21.3.6)
\begin{equation}
\label{0.22.5.5.1-fr}
  \Theta_{A'_{(K)}/K}=\Theta_{k'/k}\otimes_kK
  \tag{22.5.5.1}
\end{equation}
à un isomorphisme canonique près. En outre, on sait (21.4.7) que l'on a
$\Xi_{k'/k}=0$, autrement dit l'homomorphisme canonique (21.3.2.2)
\[
  \pi_{k'/k}:\Theta_{k'/k}\longrightarrow\Omega^1_k
\]
est injectif, si bien que l'on a une \emph{injection canonique}
\begin{equation}
\label{0.22.5.5.2-fr}
  \Theta_{A'_{(K)}/K}\longrightarrow\Omega^1_k\otimes_kK
  \tag{22.5.5.2}
\end{equation}
qui s'explicite de la façon suivante (compte tenu de (22.5.5.1))~: pour tout
$x'\in k'$, à l'élément
$d_{A'_{(K)}/K}(x'\otimes1_K)\otimes1_K$ de
$\Omega^1_{A'_{(K)}/K}\otimes_{A'_{(K)}}K$, on fait correspondre l'élément
$d_k(x'^p)\otimes1_K$ de $\Omega^1_k\otimes_kK$.
\end{env}

\begin{lemma}[22.5.6]
\label{0.22.5.6-fr}
Avec les notations de (22.5.5)~:
\begin{enumerate}
  \item[\rm(i)] Le noyau de l'homomorphisme canonique
  $\pi_{A'_{(K)}/K}$ (21.3.2.2) s'identifie par (22.5.5.2) à
  $(\Theta_{k'/k}\otimes_kK)\cap\Upsilon_{K/k}$.
  \item[\rm(ii)] Soient $A_1=A/\mathfrak m^2$,
  $A'_1=A_1\otimes_kk'=A'/\mathfrak m^2A'$~; alors l'homomorphisme
  caractéristique $\chi'_{A'_1/A_1}$ (22.4.6.2) s'identifie à la restriction de
  l'homomorphisme caractéristique
  $\chi_{A/k}:\Upsilon_{K/k}\longrightarrow V=\mathfrak m/\mathfrak m^2$
  (20.6.24) à $(\Theta_{k'/k}\otimes_kK)\cap\Upsilon_{K/k}$.
\end{enumerate}
\end{lemma}

\begin{proof}
(i) En vertu de (20.6.9), appliqué en remplaçant $A$ par le corps premier, $B$
par $k$, $C$ par $K$ et $E$ par $A_1=A/\mathfrak m^2$ (extension de $K$ par
$V=\mathfrak m/\mathfrak m^2$), il suffit (cf. (22.5.5)) de vérifier que si
$x'\in k'$ et si $z\in A_1$ est un élément dont la classe mod.
$\mathfrak m/\mathfrak m^2$ est $x'^p\in k$, alors l'image de
$d_k(x'^p)\otimes1_K$ par l'homomorphisme canonique
$\Omega^1_k\otimes_kK\longrightarrow\Omega^1_{A_1}\otimes_{A_1}K$ est
$d_{A_1}z\otimes1_K$, ce qui résulte des définitions.

(ii) Comme $A'_{(K)}=K\otimes_kk'=(A'_1)_{(K)}$ (en vertu de la relation
$K=A_1/(\mathfrak m/\mathfrak m^2)$), $\Theta_{A'_{(K)}/K}$ s'identifie à
$\Theta_{(A'_1)_{(K)}/K}$ et la vérification résulte du même calcul que dans
(i) et de la définition (22.4.5.2) de $\pi'_{(A'_1)_{(K)}/K}$.
\end{proof}

\begin{proposition}[22.5.7]
\label{0.22.5.7-fr}
Soient $k$ un corps de caractéristique $p>0$, $A$ une $k$-algèbre qui est un
anneau local régulier, $\mathfrak m$ son idéal maximal, $K=A/\mathfrak m$ son
corps résiduel~; soit $V=\mathfrak m/\mathfrak m^2$, considéré comme
$K$-espace vectoriel. Soient $k'$ une extension finie de $k$ telle que
$k'^p\subset k$, et $A'=A\otimes_kk'$~; pour que l'anneau local $A'$ soit
régulier, il faut et il suffit que la restriction de l'homomorphisme
caractéristique $\chi_{A/k}:\Upsilon_{K/k}\longrightarrow V$ (20.6.24) à
$(\Theta_{k'/k}\otimes_kK)\cap\Upsilon_{K/k}$ soit injective.
\end{proposition}

\begin{proof}
Il s'agit de prouver que la condition de l'énoncé est équivalente à
(22.5.1.1), en désignant par $\mathfrak m'$ l'idéal maximal de $A'$, par
$K'=A'/\mathfrak m'$ son corps résiduel. Avec les notations de la remarque
(22.5.2, (i)), on a $A'_1=A_1\otimes_kk'$. Or, soit
$(a_i)_{1\leqslant i\leqslant r}$ une $p$-base de $k'$ sur $k$, de sorte que
$a_i^p=b_i\in k$, et $k'$ est isomorphe au quotient de l'anneau de

\oldpage[0\textsubscript{IV-1}]{204}

polynômes $k[T_1,\ldots,T_r]$ par l'idéal engendré par les polynômes
$F_i=T_i^p-b_i$ $(1\leqslant i\leqslant r)$~; on en déduit que $A'_1$ est
isomorphe au quotient de l'anneau de polynômes $A_1[T_1,\ldots,T_r]$ par
l'idéal engendré par les $F_i$. On peut alors appliquer le critère (22.4.7), et
la condition de l'énoncé équivaut donc au fait que $\chi'_{A'_1/A_1}$ soit
injectif~; on conclut à l'aide de (22.5.6).
\end{proof}

Nous pouvons maintenant démontrer la réciproque de (19.6.5)~:

\begin{theorem}[22.5.8]
\label{0.22.5.8-fr}
Soient $k$ un corps, $p$ son exposant caractéristique, $A$ une $k$-algèbre
locale noethérienne. Les conditions suivantes sont équivalentes~:
\begin{enumerate}
  \item[a)] $A$ est une $k$-algèbre formellement lisse (pour sa topologie
  préadique).
  \item[b)] $A$ est géométriquement régulier sur $k$.
  \item[b')] Pour toute extension finie $k'$ de $k$ telle que $k'^p\subset k$,
  $A'=A\otimes_kk'$ est un anneau régulier.
\end{enumerate}
\end{theorem}

\begin{proof}
Compte tenu de (19.6.6), on peut se borner au cas $p>1$.

Le fait que a) implique b) n'est autre que (19.6.5), et b) entraîne
trivialement b'). Montrons que b') entraîne que les conditions de (22.2.2) sont
satisfaites~; c'est évident pour la première (17.1.1) puisque b') entraîne
d'abord que $A$ est régulier. D'autre part, soit $(x_\alpha)_{\alpha\in I}$ une
$p$-base de $k^{1/p}$ sur $k$~; on sait que les éléments $d_k(x_\alpha^p)$
forment une base du $k$-espace vectoriel $\Omega^1_k$ (21.4.5), donc les
éléments $d_k(x_\alpha^p)\otimes1_K$ forment une base du $K$-espace vectoriel
$\Omega^1_k\otimes_kK$. On en conclut (cf. (22.5.5)) que lorsque $k'$ parcourt
l'ensemble des sous-extensions de $k^{1/p}$, finies sur $k$, la famille des
sous-espaces $\Theta_{k'/k}\otimes_kK$ de $\Omega^1_k\otimes_kK$ est filtrante
croissante et a pour réunion $\Omega^1_k\otimes_kK$. Il résulte alors de
(22.5.7) que la condition b') entraîne que $\chi_{A/k}$ est injectif, ce qui
n'est autre que la condition (ii) de (22.2.2).
\end{proof}

\begin{corollary}[22.5.9]
\label{0.22.5.9-fr}
Soient $k$ un corps, $A$ une $k$-algèbre locale noethérienne formellement
lisse. Alors, pour tout idéal premier $\mathfrak p$ de $A$, l'anneau local
$A_{\mathfrak p}$ est une $k$-algèbre formellement lisse (pour la topologie
$\mathfrak p$-préadique).
\end{corollary}

\begin{proof}
En effet, avec les notations de (22.5.8, b')), il suffit de montrer que
l'anneau local $A_{\mathfrak p}\otimes_kk'$ est régulier~; mais comme c'est un
anneau de fractions de l'anneau local $A\otimes_kk'$ et que ce dernier est
régulier par hypothèse, il en est de même de $A_{\mathfrak p}\otimes_kk'$
(17.3.6).
\end{proof}

\begin{remark}[22.5.10]
\label{0.22.5.10-fr}
\begin{enumerate}
  \item[\rm(i)] La conclusion de (22.5.8) est en défaut lorsque, dans la
  condition b'), on se borne aux extensions \emph{monogènes} $k'=k(x)$ de $k$
  (avec $x^p\in k$). Il se peut en effet que pour toutes ces extensions $k'$,
  on ait
  $\Upsilon_{K/k}\cap(\Theta_{k'/k}\otimes_kK)=0$, bien que
  $\Upsilon_{K/k}\neq0$~; cela signifie que pour tout $x\in k^{1/p}$ tel que
  $x\notin k$, on doit avoir
  $d_k(x^p)\otimes1_K\notin\Upsilon_{K/k}$, ou encore $d_K(x^p)\neq0$,
  c'est-à-dire $x^p\notin K^p$~; autrement dit, on doit avoir
  $K^p\cap k=k^p$, bien que $K$ soit une extension \emph{inséparable} de $k$.
  Or, on construit aisément des exemples de telles extensions~: partons d'un
  corps parfait $k_0$ de caractéristique $p>0$, soient $X$, $Y$, $Z$ trois
  indéterminées et posons $k=k_0(X,Y)$ et
  $K=k(X^{1/p}+ZY^{1/p})$~; on vérifie aisément que $K$ vérifie les conditions
  précédentes, donc $\Upsilon_{K/k}\neq0$. Appliquons maintenant le lemme
  (22.2.5.2) en prenant $V=K$ et pour $h$ l'homomorphisme nul~; $A$ est alors
  un anneau de valuation discrète ayant $K$ pour corps résiduel, avec
  $\chi_{A/k}=0$, et n'est donc pas une $k$-algèbre formellement lisse par
  (22.2.2)~; toutefois $A\otimes_kk'$ est un anneau régulier pour toute
  extension \emph{monogène} $k'\subset k^{1/p}$ de $k$.

  \oldpage[0\textsubscript{IV-1}]{205}

  \item[\rm(ii)] Le corollaire (22.5.9) amène à considérer la question
  suivante~: si $A$ est une $k$-algèbre noethérienne, à quelles conditions
  l'ensemble des idéaux premiers $\mathfrak p\in\operatorname{Spec}(A)$ tels que
  $A_{\mathfrak p}$ soit une $k$-algèbre formellement lisse est-il
  \emph{ouvert}~? Nous en aborderons plus tard certains cas particuliers.
\end{enumerate}
\end{remark}

\subsection{Critère jacobien de Zariski}
\label{subsection:0.22.6-fr}

\begin{theorem}[22.6.1]
\label{0.22.6.1-fr}
\textnormal{(critère jacobien de lissité formelle).}
Soient $A$, $B$ deux anneaux topologiques,
$u:A\longrightarrow B$ un homomorphisme continu faisant de $B$ une
$A$-algèbre formellement lisse, $\mathfrak J$ un idéal de $B$ (non
nécessairement fermé), $C=B/\mathfrak J$ la $A$-algèbre topologique quotient.
Pour que $C$ soit une $A$-algèbre formellement lisse, il faut et il suffit que
l'homomorphisme canonique (cf. (20.5.11.2))
\begin{equation}
\label{0.22.6.1.1-fr}
  \delta_{C/B/A}:\mathfrak J/\mathfrak J^2
  \longrightarrow\Omega^1_{B/A}\otimes_BC
  \tag{22.6.1.1}
\end{equation}
soit \emph{formellement inversible à gauche} (cf. (19.1.5)).
\end{theorem}

\begin{proof}
En effet, dire que $B$ (resp. $C$) est une $A$-algèbre formellement lisse
signifie (19.4.4) que l'on a
$\operatorname{Exalcotop}_A(B,L)=0$ (resp.
$\operatorname{Exalcotop}_A(C,L)=0$) pour tout $B$-module (resp. $C$-module)
discret $L$ annulé par un idéal ouvert de $B$ (resp. $C$). Comme par hypothèse
$\operatorname{Exalcotop}_A(B,L)=0$ pour tout $C$-module discret $L$ annulé
par un idéal ouvert de $C$, dire que $C$ est une $A$-algèbre formellement lisse
signifie donc que l'homomorphisme canonique
$\operatorname{Exalcotop}_A(C,L)\longrightarrow
\operatorname{Exalcotop}_A(B,L)$ est injectif, et le théorème résulte donc de
(20.7.8).
\end{proof}

Rappelons (19.5.3) que lorsque $B$ et $C$ sont des $A$-algèbres formellement
lisses, $\mathfrak J/\mathfrak J^2$ est un $C$-module topologique
\emph{formellement projectif}~; en outre, les homomorphismes canoniques
$\varphi_n:\mathbf S_C^n(\mathfrak J/\mathfrak J^2)
\longrightarrow\mathfrak J^n/\mathfrak J^{n+1}$ sont des
\emph{bimorphismes formels} lorsque $B$ est supposé être un anneau
\emph{préadmissible}.
\begin{corollary}[22.6.2]
\label{0.22.6.2-fr}
Les hypothèses et notations étant celles de (22.6.1), supposons de plus que
dans $B$ le carré de tout idéal ouvert soit ouvert, et que sur $\mathfrak J$ la
topologie induite par celle de $B$ soit identique à la topologie déduite de
celle de $B$ (19.0) (on notera que ces deux conditions sont vérifiées lorsque
$B$ est noethérien et sa topologie préadique ($0_{\mathrm I}$, 7.3.2), ou si la
topologie de $B$ est la topologie $\mathfrak J$-préadique). Soit
$(\mathfrak b_\lambda)$ un système fondamental d'idéaux ouverts de $B$ et
posons $B_\lambda=B/\mathfrak b_\lambda$ pour tout $\lambda$. Alors~:
\begin{enumerate}
  \item[\rm(i)] Les conditions suivantes sont équivalentes~:
  \begin{enumerate}
    \item[a)] $C$ est une $A$-algèbre formellement lisse.
    \item[b)] Pour tout $\lambda$, l'homomorphisme de $B_\lambda$-modules
    \begin{equation}
    \label{0.22.6.2.1-fr}
      (\mathfrak J/\mathfrak J^2)\otimes_BB_\lambda
      \longrightarrow\Omega^1_{B/A}\otimes_BB_\lambda
      \tag{22.6.2.1}
    \end{equation}
    déduit de $\delta_{C/B/A}$ par tensorisation, est inversible à gauche
    (autrement dit un isomorphisme sur un facteur direct de
    $\Omega^1_{B/A}\otimes_BB_\lambda$).
  \end{enumerate}
  \item[\rm(ii)] Supposons de plus que l'anneau topologique $C$ soit
  préadmissible ($0_{\mathrm I}$, 7.1.2) et

  \oldpage[0\textsubscript{IV-1}]{206}

  soit $\mathfrak L$ un idéal de définition de $C$~; alors les conditions a)
  et b) de (i) équivalent aussi à~:
  \begin{enumerate}
    \item[c)] L'homomorphisme de $(C/\mathfrak L)$-modules
    \[
      (\mathfrak J/\mathfrak J^2)\otimes_C(C/\mathfrak L)
      \longrightarrow\Omega^1_{B/A}\otimes_B(C/\mathfrak L)
    \]
    est inversible à gauche.
  \end{enumerate}
\end{enumerate}
\end{corollary}

\begin{proof}
Les hypothèses entraînent que sur $\Omega^1_{B/A}$ la topologie est déduite de
celle de $B$ (20.4.5), donc (i) résulte de (22.6.1) et de (19.1.7). D'autre
part, rappelons (20.4.9) que puisque $B$ est une $A$-algèbre formellement
lisse, $\Omega^1_{B/A}$ est un $B$-module formellement projectif, donc
$\Omega^1_{B/A}\otimes_BB_\lambda$ est un $B_\lambda$-module projectif pour
tout $\lambda$ (19.2.4)~; l'équivalence de c) avec a) et b) lorsque $C$ est
préadmissible résulte alors de (19.1.9).
\end{proof}

En particulier~:

\begin{corollary}[22.6.3]
\label{0.22.6.3-fr}
Soient $A$ un anneau, $B$ une $A$-algèbre, $\mathfrak J$ un idéal de $B$,
$C=B/\mathfrak J$~; on suppose $A$, $B$, $C$ munis des topologies discrètes
et que $B$ est une $A$-algèbre formellement lisse. Pour que $C$ soit une
$A$-algèbre formellement lisse, il faut et il suffit que l'homomorphisme
canonique (22.6.1.1) soit inversible à gauche.
\end{corollary}

On notera que puisque toute $A$-algèbre $C$ est quotient d'une algèbre de
polynômes $A[T_\alpha]_{\alpha\in I}$, et que cette dernière est formellement
lisse (19.3.3), le critère (22.6.3) permet en principe de reconnaître si $C$
est une $A$-algèbre formellement lisse.

\begin{proposition}[22.6.4]
\label{0.22.6.4-fr}
Soient $A$ un anneau, $B$ une $A$-algèbre formellement lisse (pour les
topologies discrètes), $\mathfrak J$ un idéal de $B$, $C=B/\mathfrak J$~; on
suppose que $\mathfrak J/\mathfrak J^2$ est un $C$-module de type fini. Soient
$\mathfrak p$ un idéal premier de $C$, $k(\mathfrak p)$ le corps résiduel de
$C_{\mathfrak p}$, $\mathfrak p'$ l'image réciproque de $\mathfrak p$ dans
$B$, $\mathfrak q$ l'idéal premier de $A$ image réciproque de $\mathfrak p'$.
Les conditions suivantes sont équivalentes~:
\begin{enumerate}
  \item[a)] $C_{\mathfrak p}$ est une $A_{\mathfrak q}$-algèbre formellement
  lisse (pour les topologies discrètes).
  \item[a')] $C_{\mathfrak p}$ est une $A$-algèbre formellement lisse (pour
  les topologies discrètes).
  \item[b)] L'homomorphisme canonique
  \begin{equation}
  \label{0.22.6.4.1-fr}
    (\mathfrak J/\mathfrak J^2)\otimes_Ck(\mathfrak p)
    \longrightarrow\Omega^1_{B/A}\otimes_Bk(\mathfrak p)
    \tag{22.6.4.1}
  \end{equation}
  est injectif.
  \item[c)] Il existe $f\in C-\mathfrak p$ tel que $C_f$ soit une
  $A$-algèbre formellement lisse.
\end{enumerate}
Si de plus $B$ est noethérien, les conditions précédentes sont encore
équivalentes à~:
\begin{enumerate}
  \item[d)] $C_{\mathfrak p}$ est une $A_{\mathfrak q}$-algèbre formellement
  lisse pour la topologie $\mathfrak p$-préadique sur $C_{\mathfrak p}$ et la
  topologie discrète (ou la topologie $\mathfrak q$-préadique) sur
  $A_{\mathfrak q}$.
\end{enumerate}
\end{proposition}

\begin{proof}
Comme $A_{\mathfrak q}$ est une $A$-algèbre formellement lisse (19.3.5, (iv)
et (ii)), on sait déjà que a) et a') sont équivalentes (19.3.5, (ii)). On a
$C_{\mathfrak p}=B_{\mathfrak p'}/\mathfrak J B_{\mathfrak p'}$, et
$B_{\mathfrak p'}$ est une $A$-algèbre formellement lisse pour la topologie
discrète (19.3.5, (iv))~; en outre
$\Omega^1_{B_{\mathfrak p'}/A}=(\Omega^1_{B/A})_{\mathfrak p'}$ (20.5.9).
L'équivalence de a') et b) résulte alors de l'application de (22.6.3) à
$B_{\mathfrak p'}$ et $C_{\mathfrak p}$, en tenant compte de ce que
$\Omega^1_{B/A}$ est un $B$-module projectif (20.4.9) et (19.2.11), et que
$\mathfrak J/\mathfrak J^2$ est un $B$-module de type fini, et en utilisant
(19.1.12). L'application de (19.1.12) prouve aussi l'équivalence de b) et c).
Enfin, notons que puisque $B_{\mathfrak p'}$ est une
$A_{\mathfrak q}$-algèbre formellement lisse pour les topologies

\oldpage[0\textsubscript{IV-1}]{207}

discrètes (19.3.5, (iv)), $B_{\mathfrak p'}$, muni de la topologie
$\mathfrak p'$-préadique, est encore une $A_{\mathfrak q}$-algèbre
formellement lisse lorsque l'on prend sur $A_{\mathfrak q}$ la topologie
discrète ou la topologie $\mathfrak q$-préadique (19.3.8). Pour montrer
l'équivalence de b) et d) lorsque $B$ est noethérien, appliquons (22.6.2) à
l'anneau $A_{\mathfrak q}$ discret (ou $\mathfrak q$-préadique) et à l'anneau
$B_{\mathfrak p'}$, muni de la topologie $\mathfrak p'$-préadique~; comme
$C_{\mathfrak p}$ est préadmissible et que $\mathfrak pC_{\mathfrak p}$ est
un idéal de définition pour sa topologie, on peut invoquer l'équivalence de
c) et a) dans (22.6.2).
\end{proof}

\begin{corollary}[22.6.5]
\label{0.22.6.5-fr}
Sous les hypothèses générales de (22.6.4), l'ensemble des
$\mathfrak p\in\operatorname{Spec}(C)$ tels que $C_{\mathfrak p}$ soit une $A$-algèbre
formellement lisse (ou une $A_{\mathfrak q}$-algèbre formellement lisse, en
désignant par $\mathfrak q$ l'image réciproque de $\mathfrak p$ dans $A$)
pour les topologies discrètes, est ouvert dans $\operatorname{Spec}(C)$.
\end{corollary}

Cela résulte de la forme c) de (22.6.4). On notera que lorsque $B$ est
noethérien, on peut remplacer les topologies discrètes par les topologies
préadiques.

\begin{corollary}[22.6.6]
\label{0.22.6.6-fr}
Sous les hypothèses générales de (22.6.4), les conditions suivantes sont
équivalentes~:
\begin{enumerate}
  \item[a)] $C$ est une $A$-algèbre formellement lisse (pour les topologies
  discrètes).
  \item[b)] Pour tout $\mathfrak p\in\operatorname{Spec}(C)$ (ou seulement pour tout idéal
  maximal $\mathfrak p$ de $C$), $C_{\mathfrak p}$ est une $A$-algèbre
  formellement lisse (ou une $A_{\mathfrak q}$-algèbre formellement lisse, en
  désignant par $\mathfrak q$ l'image réciproque de $\mathfrak p$ dans $A$)
  pour les topologies discrètes.
\end{enumerate}
En outre, lorsque $B$ est noethérien, on peut remplacer dans b) les topologies
discrètes par les topologies préadiques sur $A_{\mathfrak q}$ et
$C_{\mathfrak p}$.
\end{corollary}

\begin{proof}
La dernière assertion résulte de (22.6.4). En vertu de (22.6.3), la condition
a) équivaut au fait que l'homomorphisme (22.6.1.1) est inversible à gauche~;
l'équivalence de a) et b) résulte donc de (19.1.14), compte tenu du fait que
$\Omega^1_{B/A}$ est un $B$-module projectif et $\mathfrak J/\mathfrak J^2$
un $C$-module de type fini.
\end{proof}

\begin{proposition}[22.6.7]
\label{0.22.6.7-fr}
Soient $k_0$ un corps, $k$ une extension séparable de $k_0$, $C$ une
$k$-algèbre de type fini.
\begin{enumerate}
  \item[\rm(i)] Soit $\mathfrak p$ un idéal premier de $C$. Les conditions
  suivantes sont équivalentes~:
  \begin{enumerate}
    \item[a)] $C_{\mathfrak p}$ est une $k_0$-algèbre formellement lisse pour
    la topologie discrète.
    \item[b)] $C_{\mathfrak p}$ est une $k_0$-algèbre formellement lisse pour
    la topologie $\mathfrak p$-préadique.
    \item[c)] Il existe $f\in C-\mathfrak p$ tel que $C_f$ soit une
    $k_0$-algèbre formellement lisse pour la topologie discrète.
    \item[d)] $C_{\mathfrak p}$ est un anneau géométriquement régulier
    (19.6.5) sur $k_0$.
  \end{enumerate}
  En outre l'ensemble des $\mathfrak p\in\operatorname{Spec}(C)$ ayant l'une quelconque de
  ces propriétés est ouvert dans $\operatorname{Spec}(C)$.
  \item[\rm(ii)] Les conditions suivantes sont équivalentes~:
  \begin{enumerate}
    \item[a)] $C$ est une $k_0$-algèbre formellement lisse pour la topologie
    discrète.
    \item[b)] Tout $\mathfrak p\in\operatorname{Spec}(C)$ vérifie les conditions
    équivalentes de (i).
    \item[c)] Tout idéal maximal $\mathfrak m$ de $C$ vérifie les conditions
    équivalentes de (i).
    \item[d)] Pour toute extension $k'$ de $k_0$, $C\otimes_{k_0}k'$ est un
    anneau régulier (17.3.6)~; on dit encore que $C$ est un anneau
    géométriquement régulier sur $k_0$.
  \end{enumerate}

  \oldpage[0\textsubscript{IV-1}]{208}

  \item[\rm(iii)] Soient $B=k[T_1,\ldots,T_r]$ une algèbre de polynômes,
  $\mathfrak J$ un idéal de $B$ tels que $C$ soit isomorphe à
  $B/\mathfrak J$. Soit $\mathfrak q$ un idéal premier de $B$ contenant
  $\mathfrak J$ et soit $\mathfrak p=\mathfrak q/\mathfrak J$. Les conditions
  de (i) sont alors encore équivalentes à la suivante («~critère jacobien de
  Zariski~»)~:
  \begin{enumerate}
    \item[e)] Il existe un système d'éléments
    $(f_i)_{1\leq i\leq m}$ de $\mathfrak J$, engendrant l'idéal
    $\mathfrak J_{\mathfrak q}$ dans $B_{\mathfrak q}$ et des
    $k_0$-dérivations $D_j$ de $B$ dans lui-même $(1\leq j\leq m)$ tels que
    $\det(D_j(f_i))\notin\mathfrak q$.
  \end{enumerate}
\end{enumerate}
\end{proposition}

\begin{proof}
On sait que $C$ est toujours isomorphe à une $k$-algèbre de la forme
$B/\mathfrak J$. Comme $B$ est une $k$-algèbre formellement lisse (19.3.3) et
que $k$, étant séparable sur $k_0$, est une $k_0$-algèbre formellement lisse
(19.6.1), $B$ est une $k_0$-algèbre formellement lisse (19.3.5, (ii)).
L'équivalence des conditions a), b), c) de (i) résulte donc de (22.6.4), et
celle de a), b), c) dans (ii) résulte de (22.6.6)~; l'équivalence de a) et d)
dans (i) résulte de (22.5.8)~; comme tout localisé de $C\otimes_{k_0}k'$ est
aussi un localisé de $C_{\mathfrak p}\otimes_{k_0}k'$ pour un idéal premier
$\mathfrak p\in\operatorname{Spec}(C)$ convenable, l'équivalence de d) et b) dans (ii)
résulte de l'équivalence de d) et a) dans (i). Enfin, comme
$\Omega^1_{B/k_0}$ est un $B$-module projectif, l'équivalence du critère de
Zariski e) et des autres conditions de (i) découle de (19.1.12) et de
(22.6.3), compte tenu de (20.4.8).
\end{proof}

Zariski s'intéressait en fait à un critère différentiel de régularité pour les
anneaux locaux $C_{\mathfrak p}$. Comme il revient au même de dire qu'un anneau
local noethérien contenant un corps est régulier ou est formellement lisse en
tant qu'algèbre sur son sous-corps premier (19.6.4), on obtient aussitôt un tel
critère en remplaçant $k_0$ par le sous-corps premier de $k$ dans (22.6.7)~; en
particulier, on obtient le résultat suivant, que nous retrouverons plus tard
(IV, 6.12.5) comme cas particulier de résultats plus généraux de Nagata~:

\begin{corollary}[22.6.8]
\label{0.22.6.8-fr}
\textnormal{(Zariski).}
Soit $C$ une algèbre de type fini sur un corps. Alors l'ensemble des
$\mathfrak p\in\operatorname{Spec}(C)$ tels que $C_{\mathfrak p}$ soit un anneau local
régulier est ouvert dans $\operatorname{Spec}(C)$.
\end{corollary}

\begin{remark}[22.6.9]
\label{0.22.6.9-fr}
\begin{enumerate}
  \item[\rm(i)] Les énoncés de (22.6.7) sont encore valables si au lieu de
  supposer $k$ séparable sur $k_0$, on suppose seulement qu'il est de
  \emph{multiplicité radicielle finie}, autrement dit (19.6.6) qu'il existe
  une extension radicielle finie $k'_0$ de $k_0$ telle que le corps résiduel de
  l'anneau artinien local $k\otimes_{k_0}k'_0$ soit une extension séparable
  $k'$ de $k'_0$. Il existe alors un $k'_0$-monomorphisme
  $k'\longrightarrow k\otimes_{k_0}k'_0$ qui, composé avec l'homomorphisme
  canonique $k\otimes_{k_0}k'_0\longrightarrow k'$, donne l'identité, puisque
  $k'$ est une $k'_0$-algèbre formellement lisse (19.6.1)~; on en conclut que
  $C\otimes_{k_0}k'_0$ est muni d'une structure de $k'$-algèbre de type fini,
  et comme $k'$ est séparable sur $k'_0$, on peut appliquer (22.6.7) à cette
  $k'$-algèbre~; on en conclut notre assertion en appliquant (19.4.7). Comme
  nous n'aurons pas à nous servir de cette généralisation, nous en laissons le
  détail au lecteur. Nous ignorons par contre si les résultats de (22.6.7) sont
  valables sans \emph{aucune} hypothèse sur l'extension $k$ de $k_0$.
  \item[\rm(ii)] Sous les hypothèses générales de (22.6.7), soit $(k_\alpha)$
  une famille filtrante décroissante de sous-corps de $k$ contenant $k_0$ et
  telle que
  $\bigcap_\alpha k_\alpha(k^p)=k_0(k^p)$ (où $p$ est l'exposant
  caractéristique de $k_0$). Utilisant une méthode de décompte des dimensions
  due à Nagata, on peut montrer que, dans le critère jacobien (22.6.7, (iii)),
  on peut se borner à des dérivations $D_j$ de $B$ qui soient des
  $k_\alpha$-dérivations pour un $\alpha$ convenable.

  \oldpage[0\textsubscript{IV-1}]{209}

  L'intérêt de ce résultat est qu'il y a toujours de telles familles
  $(k_\alpha)$ pour lesquelles $[k:k_\alpha]$ soit fini (21.8.6). Nous ne
  démontrerons pas ce raffinement du critère de Zariski, dont nous n'aurons pas
  à faire usage. Dans (22.7), on va donner, pour les anneaux locaux complets,
  une variante (due aussi à Nagata) du critère de Zariski, qui se démontre
  essentiellement par la même méthode (avec des difficultés techniques un peu
  plus grandes).
\end{enumerate}
\end{remark}

\subsection{Le critère jacobien de Nagata}
\label{subsection:0.22.7-fr}

\begin{env}[22.7.1]
\label{0.22.7.1-fr}
Le critère jacobien de Nagata est l'analogue du critère jacobien de Zariski,
mais pour des anneaux quotients d'anneaux de \emph{séries formelles} sur un
corps. Nous en donnerons, comme Nagata [31], deux versions, présentées ici
comme critères de lissité formelle.
\end{env}

\begin{proposition}[22.7.2]
\label{0.22.7.2-fr}
Soient $k$ un corps, $A=k[[T_1,\ldots,T_r]]$ un anneau de séries formelles sur
$k$, muni de sa structure usuelle de $k$-algèbre, $\mathfrak q$ un idéal de
$A$, $B=A/\mathfrak q$. Soit $\mathfrak p$ un idéal premier de $A$ contenant
$\mathfrak q$. On suppose qu'il existe une sous-$k$-algèbre $C'$ de
$C=A/\mathfrak p$, isomorphe à une algèbre de séries formelles
$k[[X_1,\ldots,X_s]]$, telle que $C$ soit finie sur $C'$ et que le corps des
fractions $L$ de $C$ soit une extension séparable de
$L'=k((X_1,\ldots,X_s))$ (hypothèse toujours satisfaite si $k$ est de
caractéristique 0). Alors les conditions suivantes sont équivalentes~:
\begin{enumerate}
  \item[a)] $B_{\mathfrak p}=A_{\mathfrak p}/\mathfrak qA_{\mathfrak p}$ est
  une $k$-algèbre formellement lisse pour la topologie
  $\mathfrak p$-préadique.
  \item[b)] Il existe des $k$-dérivations $D_i$ $(1\leq i\leq m)$ de $A$ dans
  lui-même, et des éléments $f_i$ $(1\leq i\leq m)$ de $\mathfrak q$, tels
  que les images des $f_i$ dans $\mathfrak qA_{\mathfrak p}$ engendrent cet
  idéal de $A_{\mathfrak p}$, et que l'on ait
  $\det(D_i(f_j))\notin\mathfrak p$.
  \item[c)] $B_{\mathfrak p}$ est un anneau régulier.
\end{enumerate}
\end{proposition}

\begin{proof}
Le corps résiduel de $B_{\mathfrak p}$ est $k(\mathfrak p)$~; comme
$k((X_1,\ldots,X_s))$ est séparable sur $k$ (21.9.6.4), $k(\mathfrak p)$ est
séparable sur $k$ par hypothèse, et l'on sait déjà que dans ces conditions les
propriétés a) et c) sont équivalentes (19.6.4). D'ailleurs
$A_{\mathfrak p}$ est un anneau régulier ((17.1.4) et (17.3.2)), et comme son
corps résiduel $k(\mathfrak p)$ est séparable sur $k$, $A_{\mathfrak p}$ est
formellement lisse sur $k$ pour la topologie $\mathfrak p$-préadique
(19.6.4). Il résulte donc de (22.6.2, (ii)) que la condition a) équivaut au
fait que l'homomorphisme de $L$-espaces vectoriels
\begin{equation}
\label{0.22.7.2.1-fr}
  j_0:(\mathfrak q/\mathfrak q^2)\otimes_AL
  \longrightarrow\Omega^1_{A_{\mathfrak p}/k}\otimes_{A_{\mathfrak p}}L
  =\widehat\Omega^1_{A/k}\otimes_AL
  \tag{22.7.2.1}
\end{equation}
soit injectif.

Considérons d'autre part l'homomorphisme composé
\begin{equation}
\label{0.22.7.2.2-fr}
  j:(\mathfrak q/\mathfrak q^2)\otimes_AL
  \xrightarrow{j_0}\Omega^1_{A/k}\otimes_AL
  \longrightarrow\widehat\Omega^1_{A/k}\otimes_AL
  \tag{22.7.2.2}
\end{equation}
et montrons que la condition b) équivaut à dire que $j$ est injectif. Notons en
effet que
\[
  (\mathfrak q/\mathfrak q^2)\otimes_AL
  =(\mathfrak qA_{\mathfrak p}/\mathfrak q^2A_{\mathfrak p})
  \otimes_{A_{\mathfrak p}}L~;
\]
la condition b) signifie que si $F_i=j(f'_i\otimes1)$, où $f'_i$ est l'image
de $f_i$ dans
$\mathfrak qA_{\mathfrak p}/\mathfrak q^2A_{\mathfrak p}$, la matrice
$(\langle F_i,D_j\otimes1\rangle)$ est inversible~; donc cela entraîne que les
$F_i$ sont linéairement indépendants, et \emph{a fortiori} qu'il en est de
même des $f'_i\otimes1$~; mais comme ces derniers engendrent
$(\mathfrak qA_{\mathfrak p}/\mathfrak q^2A_{\mathfrak p})
\otimes_{A_{\mathfrak p}}L$, on en conclut que $j$ est injectif. Inversement,
supposons $j$ injectif, et prenons les $f_i$ tels que les $f'_i\otimes1$
forment

\oldpage[0\textsubscript{IV-1}]{210}

une base de
$(\mathfrak qA_{\mathfrak p}/\mathfrak q^2A_{\mathfrak p})
\otimes_{A_{\mathfrak p}}L$~; alors les $F_i$ forment une base de l'image de
$j$~; mais on sait (21.9.3) que $\widehat\Omega^1_{A/k}$ est un $A$-module
libre de rang $r$ et les $k$-dérivations de $A$ dans lui-même engendrent son
dual~; le fait que les $F_i$ soient linéairement indépendants entraîne donc
l'existence de $k$-dérivations $D_j$ de $A$ dans lui-même telles que la matrice
$(\langle F_i,D_j\otimes1\rangle)$ soit inversible, autrement dit la condition
b).

Ces remarques montrent déjà que b) entraîne a). Pour prouver inversement que
la condition a) entraîne que $j$ est injectif, considérons le diagramme
commutatif
\begin{equation}
\label{0.22.7.2.3-fr}
\xymatrix{
  (\mathfrak q/\mathfrak q^2)\otimes_AL
    \ar[r]^j \ar[d]_\alpha &
  \widehat\Omega^1_{A/k}\otimes_AL \ar@{=}[d] \\
  (\mathfrak p/\mathfrak p^2)\otimes_AL
    \ar[r]_i &
  \widehat\Omega^1_{A/k}\otimes_AL
}
\tag{22.7.2.3}
\end{equation}
et notons le lemme suivant~:
\end{proof}

\begin{lemma}[22.7.2.4]
\label{0.22.7.2.4-fr}
Soient $R$ un anneau local régulier, $\mathfrak m$ son idéal maximal,
$K=R/\mathfrak m$ son corps résiduel. Si $\mathfrak n$ est un idéal de $R$ tel
que $R/\mathfrak n$ soit régulier, alors l'homomorphisme canonique
\begin{equation}
\label{0.22.7.2.5-fr}
  \alpha:(\mathfrak n/\mathfrak n^2)\otimes_RK
  \longrightarrow(\mathfrak m/\mathfrak m^2)\otimes_RK
  \tag{22.7.2.5}
\end{equation}
est injectif.
\end{lemma}

\begin{proof}
En effet, on sait (17.1.6) que $\mathfrak n$ est engendré par une suite
$(x_i)_{1\leq i\leq t}$ faisant partie d'un système régulier de paramètres
$(x_i)_{1\leq i\leq n}$ de $R$~; les images des $x_i$ $(1\leq i\leq t)$
forment une base de $(\mathfrak n/\mathfrak n^2)\otimes_RK$, et leurs images
par $\alpha$ font partie de la base de
$(\mathfrak m/\mathfrak m^2)\otimes_RK$ formée des images des $x_i$ pour
$1\leq i\leq n$, d'où le lemme.
\end{proof}

\begin{proof}[Suite de la démonstration de la proposition~{\rm(22.7.2)}]
Ce lemme et le diagramme (22.7.2.3) ramènent donc (en vertu de l'hypothèse a))
à prouver que $i$ est injectif. Or, dans la suite
\begin{equation}
\label{0.22.7.2.6-fr}
  \mathfrak p/\mathfrak p^2\xrightarrow{h}
  \widehat\Omega^1_{A/k}\widehat\otimes_A(A/\mathfrak p)
  \xrightarrow{g}\widehat\Omega^1_{(A/\mathfrak p)/k}\longrightarrow0
  \tag{22.7.2.6}
\end{equation}
on sait (20.7.20) que $g$ est surjectif et que l'image de $h$ est dense dans
le noyau de $g$ pour la topologie $\mathfrak m$-adique ($\mathfrak m$ étant
l'idéal maximal de $A$)~; comme $\widehat\Omega^1_{A/k}$ est un $A$-module de
type fini, tous les sous-modules du $(A/\mathfrak p)$-module
\[
  \widehat\Omega^1_{A/k}\otimes_A(A/\mathfrak p)
  =\widehat\Omega^1_{A/k}\widehat\otimes_A(A/\mathfrak p)
\]
sont fermés pour la topologie $\mathfrak m$-adique ($0_{\mathrm I}$, 7.3.5),
donc la suite (22.7.2.6) est exacte. Tensorisant par $L$, il vient la suite
exacte
\[
  (\mathfrak p/\mathfrak p^2)\otimes_AL\xrightarrow{i}
  \widehat\Omega^1_{A/k}\otimes_AL\longrightarrow
  \widehat\Omega^1_{(A/\mathfrak p)/k}\otimes_{A/\mathfrak p}L
  \longrightarrow0.
\]

Or, l'hypothèse de séparabilité faite sur $L$ entraîne, en vertu de (21.9.5),
que $\widehat\Omega^1_{(A/\mathfrak p)/k}\otimes_{A/\mathfrak p}L$ est de
rang $s$ sur $L$~; comme $\widehat\Omega^1_{A/k}\otimes_AL$ est de rang $r$
sur $L$ (21.9.3), on a $\operatorname{rg}_L(\operatorname{Im}i)=r-s$. Mais comme
$A/\mathfrak p$ est finie sur une sous-algèbre isomorphe à
$k[[T_1,\ldots,T_s]]$, on a $\dim(A/\mathfrak p)=s$ ((17.1.4) et (16.1.5))~;
donc
\[
  r-s=\dim(A)-\dim(A/\mathfrak p)=\dim(A_{\mathfrak p})
\]

\oldpage[0\textsubscript{IV-1}]{211}

par (16.5.11). Or, comme $A_{\mathfrak p}$ est régulier (17.3.2),
$\dim(A_{\mathfrak p})$ est égal au rang sur $L$ de
$\mathfrak pA_{\mathfrak p}/\mathfrak p^2A_{\mathfrak p}$ (17.1.1), donc au
rang sur $L$ de $(\mathfrak p/\mathfrak p^2)\otimes_AL$, ce qui achève de
prouver que $i$ est injectif.
\end{proof}

\begin{theorem}[22.7.3]
\label{0.22.7.3-fr}
Soient $k$ un corps, $A$ un anneau local noethérien complet de corps résiduel
$K$. On suppose que~:
\begin{enumerate}
  \item[$1^\circ$] $[k:k^p]<+\infty$ (où $p$ est l'exposant caractéristique
  de $k$)~;
  \item[$2^\circ$] $K$ est une extension finie d'une extension séparable
  $K_0$ de $k$~;
  \item[$3^\circ$] $A$ est muni d'une structure de $K_0$-algèbre formellement
  lisse (pour la topologie préadique).
\end{enumerate}
Soient $\mathfrak q$ un idéal de $A$, $B=A/\mathfrak q$, $\mathfrak p$ un
idéal premier de $A$ contenant $\mathfrak q$. Les conditions suivantes sont
équivalentes~:
\begin{enumerate}
  \item[a)] $B_{\mathfrak p}$ est une $k$-algèbre formellement lisse (pour la
  topologie $\mathfrak p$-préadique).
  \item[b)] Il existe des $k$-dérivations $D_i$ de $A$ dans lui-même
  $(1\leq i\leq m)$ et des éléments $f_i$ $(1\leq i\leq m)$ de
  $\mathfrak q$, tels que les images des $f_i$ dans
  $\mathfrak qA_{\mathfrak p}$ engendrent cet idéal de $A_{\mathfrak p}$, et
  que l'on ait $\det(D_i(f_j))\notin\mathfrak p$.
  \item[b')] Il existe une sous-extension $k'$ de $K_0$, contenant $K_0^p$,
  telle que $[K_0:k']<+\infty$, des $k'$-dérivations $D_i$ de $A$ dans
  lui-même $(1\leq i\leq m)$ et des éléments $f_i$ de $\mathfrak q$
  $(1\leq i\leq m)$, tels que les images des $f_i$ dans
  $\mathfrak qA_{\mathfrak p}$ engendrent cet idéal de $A_{\mathfrak p}$, et
  que l'on ait $\det(D_i(f_j))\notin\mathfrak p$.
\end{enumerate}
\end{theorem}

\begin{proof}
Distinguons deux cas, suivant que $p=1$ ou $p>1$.

\emph{A)} $k$ \emph{est de caractéristique} $0$. Comme $K$ est alors une
extension séparable de $K_0$, il résulte de (19.6.4) que $A$ est
$K_0$-isomorphe à un anneau de séries formelles $K[[T_1,\ldots,T_r]]$ muni de
sa structure usuelle de $K$-algèbre. Mais alors, compte tenu de (19.8.8, (ii))
appliqué à $A/\mathfrak p$, on voit que les conditions générales de (22.7.2)
sont vérifiées en y remplaçant $k$ par $K$. En outre, en vertu de (19.6.4), il
revient au même de dire que $B_{\mathfrak p}$ est une $k$-algèbre
formellement lisse ou une $K$-algèbre formellement lisse, les deux conditions
étant équivalentes au fait que $B_{\mathfrak p}$ est un anneau régulier. On
peut donc appliquer les conclusions de (22.7.2), et il est immédiat que cela
prouve l'équivalence de a), b) et b') (avec $k'=K_0$ dans b')).

\emph{B)} $k$ \emph{est de caractéristique} $p>0$. Comme $A$ est une
$K_0$-algèbre formellement lisse et que $K_0$ est séparable sur $k$, il
résulte de (19.3.5, (ii)) et (19.6.1) que $A$ est une $k$-algèbre
formellement lisse~; en vertu de (22.5.9), $A_{\mathfrak p}$ est aussi une
$k$-algèbre formellement lisse pour la topologie $\mathfrak p$-préadique. Il
résulte donc encore de (22.6.2, (ii)) que la condition a) équivaut au fait que
l'homomorphisme $j_0$ défini dans (22.7.2.1) (où $L=k(\mathfrak p)$) est
injectif. Donc (19.1.12, c)) montre donc déjà que b) implique a)~; comme par
ailleurs b') implique trivialement b), tout revient à montrer que l'hypothèse
que $j_0$ est injectif entraîne b').

Désignons d'abord par $k'$ une sous-extension quelconque de $K_0$ telle que
$[K_0:k']<+\infty$. Notons que $\widehat\Omega^1_{A/k'}$ est un $A$-module
libre de type fini. Par le raisonnement de (22.7.2), il suffit de montrer que
(pour un choix convenable de $k'$), l'homomorphisme composé
\begin{equation}
\label{0.22.7.3.1-fr}
  j':(\mathfrak q/\mathfrak q^2)\otimes_AL
  \xrightarrow{j_0}\Omega^1_{A/k}\otimes_AL
  \longrightarrow\widehat\Omega^1_{A/k'}\otimes_AL
  \tag{22.7.3.1}
\end{equation}

\oldpage[0\textsubscript{IV-1}]{212}

est injectif. Considérons encore le diagramme commutatif
\begin{equation}
\label{0.22.7.3.2-fr}
\xymatrix{
  (\mathfrak q/\mathfrak q^2)\otimes_AL
    \ar[r]^{j_0} \ar[d]_\alpha &
  \Omega^1_{A/k}\otimes_AL \ar[r] \ar@{=}[d] &
  \widehat\Omega^1_{A/k'}\otimes_AL \ar@{=}[d] \\
  (\mathfrak p/\mathfrak p^2)\otimes_AL
    \ar[r]_{i_0} &
  \Omega^1_{A/k}\otimes_AL \ar[r] &
  \widehat\Omega^1_{A/k'}\otimes_AL
}
\tag{22.7.3.2}
\end{equation}

Comme $j_0$ est injectif, on a $\operatorname{Ker}(i_0\circ\alpha)=0$. Mais comme
$A_{\mathfrak p}$ est un anneau régulier et que l'hypothèse a) implique que
$B_{\mathfrak p}=A_{\mathfrak p}/\mathfrak qA_{\mathfrak p}$ est aussi un
anneau régulier (19.6.5), le lemme (22.7.2.4) montre que $\alpha$ est injectif,
d'où $\alpha^{-1}(\operatorname{Ker} i_0)=0$. Si $i'$ est le composé des homomorphismes de la
seconde ligne de (22.7.3.2), on a de même
$\operatorname{Ker}(j')=\alpha^{-1}(\operatorname{Ker}(i'))$, et tout revient donc à montrer que, pour un
choix convenable de $k'$, on a
\begin{equation}
\label{0.22.7.3.3-fr}
  \operatorname{Ker}(i')=\operatorname{Ker}(i).
  \tag{22.7.3.3}
\end{equation}

Mais le même raisonnement que dans (22.7.2) montre que l'on a une suite exacte
\[
  (\mathfrak p/\mathfrak p^2)\otimes_AL\xrightarrow{i'}
  \widehat\Omega^1_{A/k'}\otimes_AL\longrightarrow
  \widehat\Omega^1_{(A/\mathfrak p)/k'}\otimes_{A/\mathfrak p}L
  \longrightarrow0.
\]
Or, les hypothèses faites permettent d'appliquer (21.9.8), donc il existe une
sous-extension $k'$ de $K_0$ contenant $K_0^p$, telle que
$[K_0:k']<+\infty$, et pour laquelle on a
\[
  \operatorname{rg}_L\bigl(\widehat\Omega^1_{(A/\mathfrak p)/k'}
  \otimes_{A/\mathfrak p}L\bigr)
  =\dim(A/\mathfrak p)+\operatorname{rg}_L\Upsilon_{L/k}
  +\operatorname{rg}_{K_0}\Omega^1_{K_0/k'}.
\]
Mais on a
\[
  \dim(A/\mathfrak p)=\dim(A)-\dim(A_{\mathfrak p})
\]
par (16.5.11),
\[
  \dim(A_{\mathfrak p})
  =\operatorname{rg}_L\bigl((\mathfrak p/\mathfrak p^2)\otimes_AL\bigr)
\]
puisque $A_{\mathfrak p}$ est régulier (17.1.1), et enfin
\[
  \dim(A)+\operatorname{rg}_{K_0}\Omega^1_{K_0/k'}
  =\operatorname{rg}_L\bigl(\widehat\Omega^1_{A/k'}\otimes_AL\bigr)
\]
en vertu de (21.9.2)~; on obtient donc
\begin{equation}
\label{0.22.7.3.4-fr}
  \operatorname{rg}_L(\operatorname{Ker}(i'))=\operatorname{rg}_L\Upsilon_{L/k}.
  \tag{22.7.3.4}
\end{equation}

D'autre part, le corps résiduel de $A_{\mathfrak p}$ étant formellement lisse
sur le corps premier de $k$, on a une suite exacte (20.6.22.1)
\[
  \Upsilon_{L/k}\xrightarrow{\chi_{A_{\mathfrak p}/k}}
  (\mathfrak pA_{\mathfrak p})/(\mathfrak pA_{\mathfrak p})^2
  \xrightarrow{i}\Omega^1_{A_{\mathfrak p}/k}
  \otimes_{A_{\mathfrak p}}L,
\]
et comme $A_{\mathfrak p}$ est une $k$-algèbre formellement lisse pour la
topologie $\mathfrak p$-préadique, il suit de (22.2.2) que
$\chi_{A_{\mathfrak p}/k}$ est un homomorphisme injectif, donc $\operatorname{Ker}(i)$ est
isomorphe à $\Upsilon_{L/k}$, et compte tenu de (22.7.3.4) et de ce que
$\operatorname{Ker}(i)\subset\operatorname{Ker}(i')$, cela achève de prouver (22.7.3.3) et par suite le
théorème.
\end{proof}

\oldpage[0\textsubscript{IV-1}]{213}

\begin{remark}[22.7.4]
\label{0.22.7.4-fr}
On a en fait montré (en utilisant (21.9.6)) que la condition b') est encore
équivalente aux autres conditions de (22.7.3) lorsque l'on assujettit $k'$ à
être un des corps d'une famille filtrante décroissante $(k_\alpha)$ de
sous-corps de $K_0$ contenant $K_0^p$, tels que
$[K_0:k_\alpha]<+\infty$ pour tout $\alpha$ et que
$\bigcap_\alpha k_\alpha=K_0^p$.
\end{remark}

\begin{corollary}[22.7.5]
\label{0.22.7.5-fr}
Sous les hypothèses générales de (22.7.3), l'ensemble des idéaux premiers
$\mathfrak n\in\operatorname{Spec}(B)$ tels que $B_{\mathfrak n}$ soit une
$k$-algèbre formellement lisse (pour la topologie $\mathfrak n$-préadique)
est ouvert dans $\operatorname{Spec}(B)$.
\end{corollary}

\begin{proof}
Avec les notations de (22.7.3), il suffit de voir que si $B_{\mathfrak p}$
est formellement lisse sur $k$, il en est de même de $B_{\mathfrak p'}$, pour
tout $\mathfrak p'\in V(\mathfrak q)$ assez voisin de $\mathfrak p$ dans
$\operatorname{Spec}(A)$. Or, si les $k$-dérivations $D_i$ et les éléments
$f_i$ de $\mathfrak q$ vérifient le critère b) de (22.7.3), on a aussi
$f=\det(D_if_j)\notin\mathfrak p'$ pour $\mathfrak p'\in D(f)$~; d'autre
part, il existe un $g\notin\mathfrak p$ tel que les images des $f_i$ dans
$\mathfrak qA_{\mathfrak p'}$ engendrent $\mathfrak qA_{\mathfrak p'}$ pour
$\mathfrak p'\in D(g)$ (Bourbaki, \emph{Alg. comm.}, chap.~II, §~5,
n\textsuperscript{o}~1, prop.~2), ce qui achève de prouver le corollaire.
\end{proof}

\begin{corollary}[22.7.6]
\label{0.22.7.6-fr}
\textnormal{(Nagata).} Soit $B$ un anneau local noethérien complet contenant
un corps. Alors l'ensemble des $\mathfrak n\in\operatorname{Spec}(B)$ tels
que $B_{\mathfrak n}$ soit régulier est ouvert dans $\operatorname{Spec}(B)$.
\end{corollary}

\begin{proof}
En effet, $B$ est une algèbre sur un corps premier $k$, donc (19.8.8, (ii))
de la forme $A/\mathfrak q$, où
$A=K_0[[T_1,\ldots,T_n]]$ est un anneau de séries formelles et $K_0$ une
extension de $k$~; d'autre part, dire que $B_{\mathfrak n}$ est régulier
équivaut à dire que c'est une $k$-algèbre formellement lisse pour la topologie
$\mathfrak n$-préadique (19.6.4). Toutes les conditions d'application de
(22.7.5) sont donc réalisées.
\end{proof}

\begin{remarks}[22.7.7]
\label{0.22.7.7-fr}
\begin{enumerate}
  \item[(i)] Nous verrons plus loin, avec Nagata, en utilisant (22.7.6), que
  l'on peut étendre la conclusion de (22.7.6) au cas d'un anneau local
  noethérien complet quelconque (IV, 6.12.7).

  \item[(ii)] La conclusion du théorème (22.7.3) n'est plus nécessairement
  exacte lorsqu'on ne suppose plus que $[k:k^p]$ soit fini. Prenons par
  exemple $k=\mathbf F_p(X_n)_{n\geqslant0}$,
  $A=k[[T,U]]$, $\mathfrak q$ égal à l'idéal principal $Af$, où
  \[
    f=U^p-\sum_{n=0}^{\infty}X_nT^{np}~;
  \]
  $A$ est un anneau factoriel dans lequel $f$ est un élément extrémal, car il
  est extrémal dans l'anneau de polynômes $k((T))[U]$, donc aussi dans
  l'anneau de séries formelles $k((T))[[U]]$ (Bourbaki,
  \emph{Alg. comm.}, chap.~VII, §~3), et \emph{a fortiori} dans $A$.
  Prenons $\mathfrak p=\mathfrak q$, de sorte que $B_{\mathfrak p}$ est le
  corps désigné par $E$ dans (21.9.7), qui est séparable sur $k$
  (\emph{loc. cit.}), donc une $k$-algèbre formellement lisse. Mais on a ici,
  avec les notations de (22.7.3), $k=K_0=K$, et les $k$-dérivations de $A$
  dans lui-même sont les combinaisons linéaires de $\partial/\partial T$ et
  $\partial/\partial U$~; comme on a
  $\partial f/\partial T=\partial f/\partial U=0$, les critères b) et b') de
  (22.7.3) ne sont pas vérifiés.
\end{enumerate}
\end{remarks}
